\documentclass[../main.tex]{subfiles}
\begin{document}

\chapter{GPU Simulation}
\lessoninfo{
  video={M8, lessons 1--4 (GPU Cycle Level Simulation, Parts 1 and 2; Analytical Models of GPUs; Accelerating GPU Simulation)},
  sources={Liu et al.\ \cite{liu2023photon} (Photon, assigned); Williams et al.\ \cite{williams2009roofline} (Roofline, assigned); Hong and Kim \cite{hong2009analytical} (MWP/CWP model); Huang et al.\ \cite{huang2014gpumech} (GPUMech); Karkhanis and Smith \cite{karkhanis2004firstorder} (first-order model); simulators \cite{bakhoda2009gpgpusim,khairy2020accelsim,sun2019mgpusim,gthparch2026macsim,power2015gem5gpu}; Sherwood et al.\ \cite{sherwood2002simpoint} (SimPoint); NVIDIA GEMM guide \cite{nvidia2026gemmguide}},
  textbook={Eeckhout \cite{eeckhout2010perfeval} (optional); K12 \cite{kim2012gpgpu} ch.~2 Performance Principles, §4.1.1 Execution Time Modeling, §4.2.4 Roofline Model, §4.2.6 Other GPGPU Performance Modeling Techniques; BPG \cite{nvidia2026cudabestpractices} §10.2.1 Coalesced Access to Global Memory},
  keywords={performance modeling, cycle-level simulation, execution-driven and trace-driven simulation, queue-based model, sectored cache, analytical model, Roofline model, arithmetic intensity, ridge point, CPI stack, interval analysis, MWP and CWP, sub-core, event-driven simulation, sampled simulation, SimPoint, data-driven model},
}

The mechanisms of the previous lessons (the SIMT stack and dynamic warp
formation of Lesson~6, register file caches, the warp schedulers of
Lesson~7) were all proposed and evaluated before any of them existed in
silicon.  This lesson is about the tools that make such evaluations
possible.  It first builds a cycle-level simulator of a GPU from
queues and a clock, and then adds what is specific to GPUs: warps,
divergence, coalescing and sectored caches.  It then turns to
analytical models, which predict performance from a few parameters
without simulating anything: a first-order design of a GPU, the
Roofline model, CPI stacks and interval analysis.  The last part asks
why detailed simulation is so slow and what can be done about it:
parallel and event-driven simulators, simpler models, sampling, smaller
workloads and models learned from data.

\section{Cycle-level simulation}
\label{sec:l08-cycle-level}

Before a processor is built, its designers must predict how fast it
will run real programs and how each design choice changes that
speed.\source{M8, cycle-level simulation (part~1); Eeckhout
\cite{eeckhout2010perfeval}.}  This section lists the
techniques for doing so, defines cycle-level simulation and its two
flavours, and shows how a simulator of a pipeline is organised as
queues advanced by a clock.

\subsection{Six techniques}

\begin{definition}[Performance modeling]
  \label{def:l08-performance-modeling}
  The prediction of the performance of a processor design, or of a
  program on a design, without measuring it on the finished hardware,
  is \term{performance modeling}[性能モデリング].  A program that
  imitates the behaviour of the hardware in enough detail to produce
  such a prediction is a \term{simulator}[シミュレータ].
\end{definition}

Six techniques are in common use:
\begin{enumerate}
  \item \emph{cycle-level simulation}: model what every part of the
    machine does in every cycle (this section and the next);
  \item \emph{event-driven simulation}: model only the moments at which
    something changes (\cref{sec:l08-accelerating});
  \item \emph{analytical models}: equations in a few parameters
    (\cref{sec:l08-analytical});
  \item \emph{sampling-based techniques}: simulate in detail only
    selected parts of a program (\cref{sec:l08-accelerating});
  \item \emph{data-driven, statistical or machine-learning models}:
    fit a model to measured or simulated data
    (\cref{sec:l08-accelerating});
  \item \term{FPGA emulation}[FPGAエミュレーション]: implement the
    design, or a model of it, on a field-programmable gate array, which
    runs it orders of magnitude faster than a software simulator at the
    cost of much more design effort.\margin{FPGA emulation is as
    detailed as cycle-level simulation, only faster: FPGA-based
    simulators are \enquote{fast and accurate, but they take significant
    effort to develop and use} \cite{sanchez2013zsim}.  This chapter
    does not treat them further.}
\end{enumerate}

\subsection{The cycle-level approach}

\begin{definition}[Cycle-level simulation]
  \label{def:l08-cycle-level}
  In \term{cycle-level simulation}[サイクルレベルシミュレーション] the
  simulator keeps a global clock and, in every simulated cycle, models
  the events of that cycle in each unit: which instruction is fetched,
  which is decoded, which leaves the scheduler, which request reaches
  the cache.
\end{definition}

Cycle-level simulators are the standard tool of architecture research
and are used especially early in a design, when there is no hardware
to measure.  A single clock is a simplification: a real chip runs its
cores, its memory and its on-chip network at different frequencies.
A part of a design driven by one clock is a
\term{clock domain}[クロックドメイン], and a simulator can keep one
clock per domain (core, memory, network on chip) and advance each at
its own rate.\margin{On NVIDIA GPUs the SM clock and the memory clock
are independent: \texttt{nvidia-smi -q -d CLOCK} reports
\enquote{Graphics}, \enquote{SM} and \enquote{Memory} clocks
separately.}

\subsection{Execution-driven and trace-driven simulation}

A cycle-level simulator models timing.  It must also know
\emph{which} instructions run, with which addresses and which branch
outcomes, and there are two ways to obtain them.

\begin{definition}[Execution-driven and trace-driven simulation]
  \label{def:l08-exec-trace}
  In \term{execution-driven simulation}[実行駆動シミュレーション] the
  simulator executes the program's instructions itself, computing
  their results as the simulation proceeds.  In
  \term{trace-driven simulation}[トレース駆動シミュレーション] it
  replays a \term{trace}[トレース], a record of the instructions (and
  their addresses and branch outcomes) collected beforehand from a real
  or functional execution, so that execution and timing are decoupled.
\end{definition}

An execution-driven simulator still chooses \emph{when} to execute an
instruction functionally.  It can do it when the instruction is
fetched (\emph{execute-at-fetch}), which keeps the functional model
simple, or when the timing model reaches the execute stage
(\emph{execute-at-execute}), which lets values depend on the simulated
order of events.\caution{The two choices usually give the same
result.  They differ when the simulated timing changes the value, for
example when two threads race for a lock: execute-at-fetch fixes the
winner in fetch order, execute-at-execute in the simulated execution
order.}  The two flavours trade off as follows.
\begin{description}
  \item[Execution-driven.]  No traces have to be stored, and behaviour
    that depends on the run itself is modeled faithfully: which thread
    acquires a lock first, how long a thread waits at a barrier, which
    path a data-dependent branch takes under a different schedule.
  \item[Trace-driven.]  The simulator needs no functional model, so it
    is simpler, lighter and easier to develop.  A study of caches alone
    needs only the memory addresses (a memory trace).  But the trace
    fixes the order of events as it was recorded: lock acquisition,
    barriers and anything else that depends on run-time timing cannot
    change in the simulated machine.
\end{description}
\tip{Rule of thumb: if the question is \enquote{how fast is this
design}, a trace suffices; if the design can change \emph{what} the
program does (its synchronisation or the paths its threads take), the
simulation must be execution-driven.}

\subsection{Queue-based modeling}

The pipeline of Lesson~4 is modeled most easily as stages connected by
queues (\cref{fig:l08-queue-model}).

\begin{definition}[Queue-based model]
  \label{def:l08-queue-model}
  A \term{queue-based model}[キューベースモデル] represents a pipeline
  as a set of stages connected by queues.  Each queue is a
  \term{FIFO}[FIFO] (first in, first out): instructions leave in the
  order in which they entered.  The exception is the scheduler, which
  picks any \emph{ready} instruction, so it is a pool rather than a
  FIFO.  Two parameters describe each connection: the \emph{depth} of a
  queue is the number of cycles an instruction spends in that stage,
  and the \term{pipeline width}[パイプライン幅] is the number of
  instructions that can move from one queue to the next per cycle (the
  fetch, issue or execute bandwidth).
\end{definition}

\begin{figure}[tb]
  \centering
  % Queue-based model of a pipeline (Lesson 8).  The stages fetch, decode,
% schedule, execute and write-back are connected by FIFO queues
% (FE, DE and EX queue); the scheduler is drawn as a pool from which
% ready instructions are picked, not as a FIFO.  Execute talks to the
% memory model; completion wakes up dependent instructions.
% Redraws the outline figure m8l1_queue.
% Usage: % Queue-based model of a pipeline (Lesson 8).  The stages fetch, decode,
% schedule, execute and write-back are connected by FIFO queues
% (FE, DE and EX queue); the scheduler is drawn as a pool from which
% ready instructions are picked, not as a FIFO.  Execute talks to the
% memory model; completion wakes up dependent instructions.
% Redraws the outline figure m8l1_queue.
% Usage: % Queue-based model of a pipeline (Lesson 8).  The stages fetch, decode,
% schedule, execute and write-back are connected by FIFO queues
% (FE, DE and EX queue); the scheduler is drawn as a pool from which
% ready instructions are picked, not as a FIFO.  Execute talks to the
% memory model; completion wakes up dependent instructions.
% Redraws the outline figure m8l1_queue.
% Usage: \input{figures/l08-queue-model}
% Width: body (120 mm); the drawing is about 116 mm wide.
\begin{tikzpicture}[x=1mm, y=1mm, font=\scriptsize\sffamily,
    >={Stealth[length=1.5mm]},
    st/.style={draw, fill=ln-accent!15, minimum width=14mm,
               minimum height=8mm, align=center, inner sep=1pt},
    q/.style={draw, fill=ln-box, minimum width=12mm, minimum height=4mm,
              inner sep=0pt}]
  % \fifo{x}{y}{name}{label}: a queue drawn as four slots
  \newcommand{\fifo}[4]{%
    \foreach \i in {0,...,3} {
      \draw[fill=ln-box] (#1+\i*3,#2-2) rectangle ++(3,4);
    }
    \coordinate (#3-in) at (#1,#2);
    \coordinate (#3-out) at (#1+12,#2);
    \node[above, inner sep=1pt] at (#1+6,#2+2) {#4};}
  \node[st] (fe) at (7,20) {\iflangja{フェッチ}{fetch}};
  \fifo{18}{20}{feq}{FE queue}
  \node[st] (de) at (41,20) {\iflangja{デコード}{decode}};
  \fifo{52}{20}{deq}{DE queue}
  % scheduler: a pool, not a FIFO
  \node[draw, fill=ln-caution!12, rounded corners=2mm, minimum width=18mm,
        minimum height=10mm, align=center, inner sep=1pt] (sch) at (79,20)
    {\iflangja{スケジューラ}{scheduler}\\[-0.3mm]%
     \iflangja{（FIFO ではない）}{(not a FIFO)}};
  \node[st] (ex) at (103,20) {\iflangja{実行}{execute}};
  \node[st, fill=ln-tip!15] (mem) at (103,4)
    {\iflangja{キャッシュ・メモリ}{caches, memory}};
  \fifo{70}{-2}{exq}{}
  \node[below, inner sep=1pt] at (76,-4) {EX queue};
  \node[st] (wb) at (49,-2) {\iflangja{ライトバック}{write-back}};
  \draw[->] (fe) -- (feq-in);
  \draw[->] (feq-out) -- (de);
  \draw[->] (de) -- (deq-in);
  \draw[->] (deq-out) -- (sch);
  \draw[->] (sch) -- (ex);
  \draw[<->] (ex) -- (mem);
  \draw[->] (ex.east) -- ++(4,0) |- (exq-out);
  \draw[->] (exq-in) -- (wb);
  % wake-up of dependants
  \draw[->, dashed, ln-caution] (wb.north) -- ++(0,5) -| (sch.south)
    node[pos=0.25, above, text=ln-caution, inner sep=1.5pt]
    {\iflangja{完了で依存命令を起こす}{completion wakes dependants}};
  % legend: depth and width
  \draw[|-|, ln-accent] (18,27) -- (30,27)
    node[midway, above, text=ln-accent, inner sep=1.5pt]
    {\iflangja{深さ = ステージの遅延}{depth = stage latency}};
  \draw[->, thick, ln-accent] (45,27) -- (53,27)
    node[right, text=ln-accent, inner sep=1.5pt]
    {\iflangja{幅 = 1サイクルに移動できる命令数}{width = instructions moved per cycle}};
\end{tikzpicture}

% Width: body (120 mm); the drawing is about 116 mm wide.
\begin{tikzpicture}[x=1mm, y=1mm, font=\scriptsize\sffamily,
    >={Stealth[length=1.5mm]},
    st/.style={draw, fill=ln-accent!15, minimum width=14mm,
               minimum height=8mm, align=center, inner sep=1pt},
    q/.style={draw, fill=ln-box, minimum width=12mm, minimum height=4mm,
              inner sep=0pt}]
  % \fifo{x}{y}{name}{label}: a queue drawn as four slots
  \newcommand{\fifo}[4]{%
    \foreach \i in {0,...,3} {
      \draw[fill=ln-box] (#1+\i*3,#2-2) rectangle ++(3,4);
    }
    \coordinate (#3-in) at (#1,#2);
    \coordinate (#3-out) at (#1+12,#2);
    \node[above, inner sep=1pt] at (#1+6,#2+2) {#4};}
  \node[st] (fe) at (7,20) {\iflangja{フェッチ}{fetch}};
  \fifo{18}{20}{feq}{FE queue}
  \node[st] (de) at (41,20) {\iflangja{デコード}{decode}};
  \fifo{52}{20}{deq}{DE queue}
  % scheduler: a pool, not a FIFO
  \node[draw, fill=ln-caution!12, rounded corners=2mm, minimum width=18mm,
        minimum height=10mm, align=center, inner sep=1pt] (sch) at (79,20)
    {\iflangja{スケジューラ}{scheduler}\\[-0.3mm]%
     \iflangja{（FIFO ではない）}{(not a FIFO)}};
  \node[st] (ex) at (103,20) {\iflangja{実行}{execute}};
  \node[st, fill=ln-tip!15] (mem) at (103,4)
    {\iflangja{キャッシュ・メモリ}{caches, memory}};
  \fifo{70}{-2}{exq}{}
  \node[below, inner sep=1pt] at (76,-4) {EX queue};
  \node[st] (wb) at (49,-2) {\iflangja{ライトバック}{write-back}};
  \draw[->] (fe) -- (feq-in);
  \draw[->] (feq-out) -- (de);
  \draw[->] (de) -- (deq-in);
  \draw[->] (deq-out) -- (sch);
  \draw[->] (sch) -- (ex);
  \draw[<->] (ex) -- (mem);
  \draw[->] (ex.east) -- ++(4,0) |- (exq-out);
  \draw[->] (exq-in) -- (wb);
  % wake-up of dependants
  \draw[->, dashed, ln-caution] (wb.north) -- ++(0,5) -| (sch.south)
    node[pos=0.25, above, text=ln-caution, inner sep=1.5pt]
    {\iflangja{完了で依存命令を起こす}{completion wakes dependants}};
  % legend: depth and width
  \draw[|-|, ln-accent] (18,27) -- (30,27)
    node[midway, above, text=ln-accent, inner sep=1.5pt]
    {\iflangja{深さ = ステージの遅延}{depth = stage latency}};
  \draw[->, thick, ln-accent] (45,27) -- (53,27)
    node[right, text=ln-accent, inner sep=1.5pt]
    {\iflangja{幅 = 1サイクルに移動できる命令数}{width = instructions moved per cycle}};
\end{tikzpicture}

% Width: body (120 mm); the drawing is about 116 mm wide.
\begin{tikzpicture}[x=1mm, y=1mm, font=\scriptsize\sffamily,
    >={Stealth[length=1.5mm]},
    st/.style={draw, fill=ln-accent!15, minimum width=14mm,
               minimum height=8mm, align=center, inner sep=1pt},
    q/.style={draw, fill=ln-box, minimum width=12mm, minimum height=4mm,
              inner sep=0pt}]
  % \fifo{x}{y}{name}{label}: a queue drawn as four slots
  \newcommand{\fifo}[4]{%
    \foreach \i in {0,...,3} {
      \draw[fill=ln-box] (#1+\i*3,#2-2) rectangle ++(3,4);
    }
    \coordinate (#3-in) at (#1,#2);
    \coordinate (#3-out) at (#1+12,#2);
    \node[above, inner sep=1pt] at (#1+6,#2+2) {#4};}
  \node[st] (fe) at (7,20) {\iflangja{フェッチ}{fetch}};
  \fifo{18}{20}{feq}{FE queue}
  \node[st] (de) at (41,20) {\iflangja{デコード}{decode}};
  \fifo{52}{20}{deq}{DE queue}
  % scheduler: a pool, not a FIFO
  \node[draw, fill=ln-caution!12, rounded corners=2mm, minimum width=18mm,
        minimum height=10mm, align=center, inner sep=1pt] (sch) at (79,20)
    {\iflangja{スケジューラ}{scheduler}\\[-0.3mm]%
     \iflangja{（FIFO ではない）}{(not a FIFO)}};
  \node[st] (ex) at (103,20) {\iflangja{実行}{execute}};
  \node[st, fill=ln-tip!15] (mem) at (103,4)
    {\iflangja{キャッシュ・メモリ}{caches, memory}};
  \fifo{70}{-2}{exq}{}
  \node[below, inner sep=1pt] at (76,-4) {EX queue};
  \node[st] (wb) at (49,-2) {\iflangja{ライトバック}{write-back}};
  \draw[->] (fe) -- (feq-in);
  \draw[->] (feq-out) -- (de);
  \draw[->] (de) -- (deq-in);
  \draw[->] (deq-out) -- (sch);
  \draw[->] (sch) -- (ex);
  \draw[<->] (ex) -- (mem);
  \draw[->] (ex.east) -- ++(4,0) |- (exq-out);
  \draw[->] (exq-in) -- (wb);
  % wake-up of dependants
  \draw[->, dashed, ln-caution] (wb.north) -- ++(0,5) -| (sch.south)
    node[pos=0.25, above, text=ln-caution, inner sep=1.5pt]
    {\iflangja{完了で依存命令を起こす}{completion wakes dependants}};
  % legend: depth and width
  \draw[|-|, ln-accent] (18,27) -- (30,27)
    node[midway, above, text=ln-accent, inner sep=1.5pt]
    {\iflangja{深さ = ステージの遅延}{depth = stage latency}};
  \draw[->, thick, ln-accent] (45,27) -- (53,27)
    node[right, text=ln-accent, inner sep=1.5pt]
    {\iflangja{幅 = 1サイクルに移動できる命令数}{width = instructions moved per cycle}};
\end{tikzpicture}

  \caption{A queue-based model of a pipeline.  Fetch, decode,
    schedule, execute and write-back are connected by FIFO queues; the
    scheduler is a pool from which ready instructions are picked.  The
    execute stage obtains memory latencies from a model of the caches
    and memory, and a completing instruction wakes up its dependants.}
  \label{fig:l08-queue-model}
\end{figure}

Three things give the model its timing.  First, dependences: when an
instruction completes, its result must be broadcast to the
instructions waiting for it, which then become ready, exactly as the
scoreboard of Lesson~4 does in hardware.  Second, the latency of memory
instructions comes from a model of the caches and memory, which
decides hit or miss at each level.  Third, the latency of other
instructions is a parameter, for example \num{1}~cycle for an integer
add and \num{3}~cycles for a multiply, taken from the literature,
from a circuit model such as CACTI, or from the register-transfer-level
(RTL) description of the design; the simulator itself simulates no
circuits.\margin{CACTI is a widely used model of the access time,
area and energy of SRAM arrays such as caches and register files,
from their size and organisation.}  Execution latencies that do not
depend on the run are therefore a lookup table.

\subsection{Structure of a simulator}

A cycle-level simulator is a loop that advances the clock by one cycle
per iteration and asks every component to do one cycle's work.  The
components form a hierarchy: the simulator ticks the cores, the memory
system and the network on chip, and each core ticks its pipeline
stages.\margin{MacSim, a CPU and GPU simulator from Georgia Tech
\cite{gthparch2026macsim}, is organised this way; the listing is our
own sketch, not its code.}

\begin{lstlisting}[language=C++, numbers=none]
// Top level: one iteration = one simulated cycle.
while (!sim.all_done()) {
  ++sim.cycle;
  for (Core &c : sim.cores) c.tick();
  sim.memory.tick();
  sim.noc.tick();
}

// One core: stages in REVERSE pipeline order.
void Core::tick() {
  writeback(); execute(); schedule();
  decode(); fetch();
}
\end{lstlisting}

The stages are called from the last to the first.  When
\texttt{execute()} runs, it consumes the instructions that
\texttt{schedule()} produced in the \emph{previous} cycle, before
\texttt{schedule()} overwrites them with this cycle's output.  Calling
them in pipeline order would let an instruction pass through all
five stages in a single cycle.\tip{Reverse order is the software
equivalent of the pipeline registers between stages: each stage reads
what the stage before it wrote one cycle earlier.}

Each simulated instruction is represented by a small record (an
\emph{op}) that carries it through the stages, together with its
timing state.  The key field is the cycle at which it completes:
\begin{equation}
  \textit{done\_cycle} = \textit{schedule\_cycle} + \textit{latency}.
  \label{eq:l08-done-cycle}
\end{equation}
For an instruction of fixed latency the scheduler can compute it at
the moment of issue.  The latency of a load is not fixed, since it
depends on whether the access hits in the L1 cache, the L2 cache or
neither; it becomes known only after the access has been simulated
in the cache hierarchy, and is then added to the issue cycle.

\begin{example}
  \label{ex:l08-done-cycle}
  With latencies of \num{1}~cycle for \texttt{add}, \num{3} for
  \texttt{mul}, and \num{4}, \num{30} or \num{200}~cycles for a load
  that hits in L1, hits in L2 or misses in both (our own numbers), a
  simulator schedules four instructions as in
  \cref{tab:l08-done-cycle}.  The \texttt{add} depends on the load, so
  it becomes ready only at the load's \textit{done\_cycle}.
\end{example}

\begin{table}[tb]
  \centering
  \caption{Completion times computed by a simulator for four
    instructions (our own latencies).  The load's latency is known only
    after the cache model has decided where it hits.}
  \label{tab:l08-done-cycle}
  \begin{tabular}{@{}l c c c >{\raggedright\arraybackslash}p{40mm}@{}}
    \toprule
    instruction & schedule & latency & done & latency from \\
    \midrule
    \texttt{mul r1, r2, r3} & 10 & 3 & 13 & table \\
    \texttt{ld r4, [r5]} & 11 & 30 & 41 & cache model: L1 miss, L2 hit \\
    \texttt{mul r6, r1, r1} & 13 & 3 & 16 & table; waits for \texttt{r1} \\
    \texttt{add r7, r4, 1} & 41 & 1 & 42 & table; waits for \texttt{r4} \\
    \bottomrule
  \end{tabular}
\end{table}

The scheduler is the most detailed part of the model.  Each cycle it
selects instructions that are ready, subject to the resources they
need: a free floating-point unit, an entry in the load or store queue,
a free execution or cache port.  An instruction that occupies a unit for
several cycles keeps the resource for all of them, and the scheduler
must not issue a second instruction to it in the meantime.

\subsection{Modeling a GPU}

A cycle-level GPU simulator has the same structure, with one change
of unit: it models \emph{warps}, not threads, because the SIMT hardware
fetches, decodes, schedules and executes one warp instruction for all
its threads together (Lesson~4).  The scheduler considers the
instruction at the head of each warp; instructions of one warp issue in
order, those of different warps in any order.\margin{A GPU simulator
has far more to model per cycle than a CPU simulator: every SM can
hold up to 64 resident warps (CPG Table~21 \cite{nvidia2024cudaguide}),
and a GPU has tens of SMs.}  Three things have no CPU
counterpart: divergent warps (\cref{sec:l08-divergence-memory}), the
hierarchy of warps, thread blocks and kernels, and the warp scheduling
policies of Lesson~7.

The hierarchy also decides when a simulation ends.  A whole thread
block is assigned to one SM (Lesson~3), and the simulator tracks every
thread until it executes its last instruction.\margin{How many blocks
an SM holds at once is itself a modeled limit: at most 32 resident
blocks per SM on compute capabilities 7.0, 8.0 and 9.0, 16 on 7.5 and
8.6 (CPG Table~21 \cite{nvidia2024cudaguide}).}  A block finishes when
all of its threads have finished, which frees its SM resources for the
next block; a kernel finishes when all of its blocks have finished; and
the application finishes when all of its kernels have.

\begin{keypoint}
  A cycle-level simulator advances a global clock (one per clock
  domain) and models every unit in every cycle.  It obtains the
  instruction stream by executing the program or by replaying a trace.
  The pipeline is a set of FIFO queues whose depths are stage
  latencies and whose transfer rates are the pipeline widths, plus a
  scheduler that picks ready instructions under resource constraints.
  A GPU simulator models warps instead of threads.
\end{keypoint}

\section{Divergence and memory in GPU simulators}
\label{sec:l08-divergence-memory}

The modeling unit of a GPU simulator is the warp: an instruction-cache
access and a program counter belong to a warp, not to a
thread.\source{M8, cycle-level simulation (part~2); BPG
§10.2.1 Coalesced Access to Global Memory
\cite{nvidia2026cudabestpractices}.}  Not every thread of a warp takes
part in every instruction, so the simulator also keeps the active mask
of Lesson~4, both to compute results and to count the resources an
instruction really uses.  This section shows how a simulator handles
the two GPU-specific effects of the previous lessons, divergence and
coalescing, and then lists the open-source GPU simulators.

\subsection{Divergent branches}

How a simulator handles a divergent branch (Lesson~6) depends on how
it obtains its instructions.
\begin{description}
  \item[Execution-driven.]  The simulator computes each thread's branch
    outcome and maintains the SIMT stack itself, exactly as the
    hardware would.  Because the stack is part of the model, a
    different mechanism (dynamic warp formation, another reconvergence
    rule) can be modeled by replacing it.
  \item[Trace-driven.]  The trace must record every path the warp
    executed, with the active mask of each basic block, for example
    BB1 \texttt{1111}, BB2 \texttt{1100}, BB3 \texttt{0011} and BB4
    \texttt{1111} for the branch of Lesson~6.  The simulator replays
    these blocks in the recorded order.
\end{description}
\caution{A trace records the divergence behaviour of the machine on
which it was collected.  Simulating a different divergence-handling
mechanism from such a trace is difficult, because the trace contains
neither the branch conditions nor the order the new mechanism would
choose.}%
A trace-driven simulator is therefore poorly suited to studying
divergence handling itself; it is adequate when divergence is only one
of the effects that shape the timing.

\subsection{Coalescing}

The number of memory transactions a warp instruction generates, not
the number of threads, drives everything downstream: cache accesses,
misses, DRAM traffic and bandwidth.  Modeling memory coalescing
(Lesson~4) is therefore one of the most important parts of a GPU
simulator.  Requests are merged at the granularity of the first cache
they reach, typically one cache line.  In the toy configuration of this
section, with \SI{64}{\byte} lines and eight threads per warp, a load
\texttt{R1} whose threads read consecutive \SI{4}{\byte} words is
merged into one \SI{64}{\byte} request, while a load \texttt{R2} whose
threads read words \SI{128}{\byte} apart produces eight requests, one
per thread, each for a full line.

A trace can represent the memory accesses of a warp instruction in
two ways.  It can list the address of every thread (\texttt{0x00},
\texttt{0x04}, \texttt{0x08}, \dots), leaving the coalescing to the
simulator, or it can store the result of coalescing, one base address
and a size per transaction.  The second form is much smaller, but it
fixes the transaction size of the machine on which the trace was
collected.\caution{A pre-coalesced trace is tied to one
microarchitecture: to study a cache with another line or transaction
size, new traces must be collected.}

\subsection{Cache hierarchy and sectored caches}

After coalescing, the requests enter a cache hierarchy that is modeled
much as for a CPU: the address is translated by a TLB (with a page
table walk on a miss, Lesson~7), then the L1 and L2 caches are
accessed, depending on the microarchitecture being modeled.  What
differs is the importance of coalescing in front of it, and one
feature of the caches themselves.

\begin{definition}[Sectored cache]
  \label{def:l08-sectored-cache}
  A \term{sectored cache}[セクタキャッシュ] divides each cache block
  into several sectors, each with its own valid bit, all sharing the
  block's single tag.  A \term{sector}[セクタ] is the unit that is
  filled from memory on a miss: the cache brings in only the sector
  that holds the requested data rather than the whole block
  (\cref{fig:l08-sector-cache}).
\end{definition}

\begin{figure}[tb]
  \centering
  % Coalescing and sectored caches in a simulator (Lesson 8).  Top: one
% 64-byte cache block of a sectored cache, four 16-byte sectors sharing
% one tag, each with its own valid bit.  Bottom: the uncoalesced load R2
% (eight threads, addresses 128 B apart) misses in eight blocks; a
% conventional cache fills eight whole blocks (512 B), a sectored cache
% only the eight sectors that hold the requested words (128 B).
% Redraws the outline figures m8l2_memory_coalescing and
% m8l2_sector_cache (toy sizes of the lecture: 64 B block, 16 B sector;
% sector address typo corrected).
% Usage: % Coalescing and sectored caches in a simulator (Lesson 8).  Top: one
% 64-byte cache block of a sectored cache, four 16-byte sectors sharing
% one tag, each with its own valid bit.  Bottom: the uncoalesced load R2
% (eight threads, addresses 128 B apart) misses in eight blocks; a
% conventional cache fills eight whole blocks (512 B), a sectored cache
% only the eight sectors that hold the requested words (128 B).
% Redraws the outline figures m8l2_memory_coalescing and
% m8l2_sector_cache (toy sizes of the lecture: 64 B block, 16 B sector;
% sector address typo corrected).
% Usage: % Coalescing and sectored caches in a simulator (Lesson 8).  Top: one
% 64-byte cache block of a sectored cache, four 16-byte sectors sharing
% one tag, each with its own valid bit.  Bottom: the uncoalesced load R2
% (eight threads, addresses 128 B apart) misses in eight blocks; a
% conventional cache fills eight whole blocks (512 B), a sectored cache
% only the eight sectors that hold the requested words (128 B).
% Redraws the outline figures m8l2_memory_coalescing and
% m8l2_sector_cache (toy sizes of the lecture: 64 B block, 16 B sector;
% sector address typo corrected).
% Usage: \input{figures/l08-sector-cache}
% Width: body (120 mm); the drawing is about 114 mm wide.
\begin{tikzpicture}[x=1mm, y=1mm, font=\scriptsize\sffamily]
  % one block: tag + 4 sectors with valid bits
  \node[anchor=west, font=\scriptsize\sffamily\bfseries] at (0,51)
    {\iflangja{(a) セクタキャッシュのブロック（64\,B）}{(a) one block of a sectored cache (64\,B)}};
  \draw[fill=ln-accent!15] (0,40) rectangle (16,46) node[midway] {\iflangja{タグ}{tag}};
  \foreach \i in {0,...,3} {
    \draw[fill=ln-box] (17+\i*24,40) rectangle ++(4,6) node[midway] {V};
    \draw[fill=ln-tip!15] (21+\i*24,40) rectangle ++(19,6)
      node[midway] {\iflangja{セクタ}{sector} \i: 16\,B};
  }
  \node[anchor=west, text=ln-muted] at (0,36.5)
    {\iflangja{4つのセクタが1つのタグを共有し，セクタごとに有効ビットを持つ}{four sectors share one tag; each sector has its own valid bit}};
  % R2: eight misses
  \node[anchor=west, font=\scriptsize\sffamily\bfseries] at (0,31)
    {\iflangja{(b) R2 のミス（アドレス 0, 128, \dots, 896）}{(b) misses of R2 (addresses 0, 128, \dots, 896)}};
  \node[anchor=west] at (0,21) {\iflangja{従来型}{conventional}};
  \node[anchor=west] at (0,13) {\iflangja{セクタ型}{sectored}};
  \foreach \k in {0,...,7} {
    \pgfmathsetmacro{\x}{20+\k*11.5}
    \node[text=ln-muted] at (\x+5,26) {\pgfmathparse{int(\k*128)}\pgfmathresult};
    \foreach \s in {0,...,3} {
      \draw[fill=ln-caution!25] (\x+\s*2.5,18) rectangle ++(2.5,6);
      \ifnum\s=0
        \draw[fill=ln-caution!25] (\x+\s*2.5,10) rectangle ++(2.5,6);
      \else
        \draw[fill=white] (\x+\s*2.5,10) rectangle ++(2.5,6);
      \fi
    }
  }
  \node[anchor=west] at (20,5.5)
    {\iflangja{従来型: $8 \times 64 = 512$\,B を転送}{conventional: $8 \times 64 = 512$\,B fetched}};
  \node[anchor=west] at (68,5.5)
    {\iflangja{セクタ型: $8 \times 16 = 128$\,B を転送}{sectored: $8 \times 16 = 128$\,B fetched}};
  \node[anchor=west, text=ln-muted] at (20,1)
    {\iflangja{使われるのは各ブロックの 4\,B（合計 32\,B）}{each block supplies one 4\,B word (32\,B in total)}};
\end{tikzpicture}

% Width: body (120 mm); the drawing is about 114 mm wide.
\begin{tikzpicture}[x=1mm, y=1mm, font=\scriptsize\sffamily]
  % one block: tag + 4 sectors with valid bits
  \node[anchor=west, font=\scriptsize\sffamily\bfseries] at (0,51)
    {\iflangja{(a) セクタキャッシュのブロック（64\,B）}{(a) one block of a sectored cache (64\,B)}};
  \draw[fill=ln-accent!15] (0,40) rectangle (16,46) node[midway] {\iflangja{タグ}{tag}};
  \foreach \i in {0,...,3} {
    \draw[fill=ln-box] (17+\i*24,40) rectangle ++(4,6) node[midway] {V};
    \draw[fill=ln-tip!15] (21+\i*24,40) rectangle ++(19,6)
      node[midway] {\iflangja{セクタ}{sector} \i: 16\,B};
  }
  \node[anchor=west, text=ln-muted] at (0,36.5)
    {\iflangja{4つのセクタが1つのタグを共有し，セクタごとに有効ビットを持つ}{four sectors share one tag; each sector has its own valid bit}};
  % R2: eight misses
  \node[anchor=west, font=\scriptsize\sffamily\bfseries] at (0,31)
    {\iflangja{(b) R2 のミス（アドレス 0, 128, \dots, 896）}{(b) misses of R2 (addresses 0, 128, \dots, 896)}};
  \node[anchor=west] at (0,21) {\iflangja{従来型}{conventional}};
  \node[anchor=west] at (0,13) {\iflangja{セクタ型}{sectored}};
  \foreach \k in {0,...,7} {
    \pgfmathsetmacro{\x}{20+\k*11.5}
    \node[text=ln-muted] at (\x+5,26) {\pgfmathparse{int(\k*128)}\pgfmathresult};
    \foreach \s in {0,...,3} {
      \draw[fill=ln-caution!25] (\x+\s*2.5,18) rectangle ++(2.5,6);
      \ifnum\s=0
        \draw[fill=ln-caution!25] (\x+\s*2.5,10) rectangle ++(2.5,6);
      \else
        \draw[fill=white] (\x+\s*2.5,10) rectangle ++(2.5,6);
      \fi
    }
  }
  \node[anchor=west] at (20,5.5)
    {\iflangja{従来型: $8 \times 64 = 512$\,B を転送}{conventional: $8 \times 64 = 512$\,B fetched}};
  \node[anchor=west] at (68,5.5)
    {\iflangja{セクタ型: $8 \times 16 = 128$\,B を転送}{sectored: $8 \times 16 = 128$\,B fetched}};
  \node[anchor=west, text=ln-muted] at (20,1)
    {\iflangja{使われるのは各ブロックの 4\,B（合計 32\,B）}{each block supplies one 4\,B word (32\,B in total)}};
\end{tikzpicture}

% Width: body (120 mm); the drawing is about 114 mm wide.
\begin{tikzpicture}[x=1mm, y=1mm, font=\scriptsize\sffamily]
  % one block: tag + 4 sectors with valid bits
  \node[anchor=west, font=\scriptsize\sffamily\bfseries] at (0,51)
    {\iflangja{(a) セクタキャッシュのブロック（64\,B）}{(a) one block of a sectored cache (64\,B)}};
  \draw[fill=ln-accent!15] (0,40) rectangle (16,46) node[midway] {\iflangja{タグ}{tag}};
  \foreach \i in {0,...,3} {
    \draw[fill=ln-box] (17+\i*24,40) rectangle ++(4,6) node[midway] {V};
    \draw[fill=ln-tip!15] (21+\i*24,40) rectangle ++(19,6)
      node[midway] {\iflangja{セクタ}{sector} \i: 16\,B};
  }
  \node[anchor=west, text=ln-muted] at (0,36.5)
    {\iflangja{4つのセクタが1つのタグを共有し，セクタごとに有効ビットを持つ}{four sectors share one tag; each sector has its own valid bit}};
  % R2: eight misses
  \node[anchor=west, font=\scriptsize\sffamily\bfseries] at (0,31)
    {\iflangja{(b) R2 のミス（アドレス 0, 128, \dots, 896）}{(b) misses of R2 (addresses 0, 128, \dots, 896)}};
  \node[anchor=west] at (0,21) {\iflangja{従来型}{conventional}};
  \node[anchor=west] at (0,13) {\iflangja{セクタ型}{sectored}};
  \foreach \k in {0,...,7} {
    \pgfmathsetmacro{\x}{20+\k*11.5}
    \node[text=ln-muted] at (\x+5,26) {\pgfmathparse{int(\k*128)}\pgfmathresult};
    \foreach \s in {0,...,3} {
      \draw[fill=ln-caution!25] (\x+\s*2.5,18) rectangle ++(2.5,6);
      \ifnum\s=0
        \draw[fill=ln-caution!25] (\x+\s*2.5,10) rectangle ++(2.5,6);
      \else
        \draw[fill=white] (\x+\s*2.5,10) rectangle ++(2.5,6);
      \fi
    }
  }
  \node[anchor=west] at (20,5.5)
    {\iflangja{従来型: $8 \times 64 = 512$\,B を転送}{conventional: $8 \times 64 = 512$\,B fetched}};
  \node[anchor=west] at (68,5.5)
    {\iflangja{セクタ型: $8 \times 16 = 128$\,B を転送}{sectored: $8 \times 16 = 128$\,B fetched}};
  \node[anchor=west, text=ln-muted] at (20,1)
    {\iflangja{使われるのは各ブロックの 4\,B（合計 32\,B）}{each block supplies one 4\,B word (32\,B in total)}};
\end{tikzpicture}

  \caption{A sectored cache (toy sizes: \SI{64}{\byte} block,
    \SI{16}{\byte} sectors).  (a) The sectors share one tag.  (b) The
    uncoalesced load \texttt{R2} misses in eight blocks; a conventional
    cache fills eight whole blocks, a sectored cache only one sector
    in each.}
  \label{fig:l08-sector-cache}
\end{figure}

For the coalesced load \texttt{R1} (eight consecutive \SI{4}{\byte}
words, \SI{32}{\byte} at addresses 0 to 31) the conventional cache fills
the whole \SI{64}{\byte} block, while the sectored cache fills only
sectors 0 and~1, the \SI{32}{\byte} that hold the data.  The saving is
a factor of two, and every fetched byte is used.  For the uncoalesced load \texttt{R2} the conventional cache
fetches $8 \times 64 = \SI{512}{\byte}$ to deliver $8 \times 4 =
\SI{32}{\byte}$ of useful data, the sectored cache
$8 \times 16 = \SI{128}{\byte}$.  The benefit of sectoring is
therefore lower memory bandwidth for scattered accesses, and sharing
one tag among several sectors keeps the tag storage of a large block.
The cost is less implicit prefetching: a block that is filled sector by
sector no longer brings in the neighbouring data that spatial locality
would have used soon after.\caution{Lecture vs guide: the
\SI{16}{\byte} sector in a \SI{64}{\byte} block is a toy size.  NVIDIA
GPUs use \SI{32}{\byte} sectors in \SI{128}{\byte} lines, and since
compute capability 6.0 a warp's accesses cost \enquote{the number of
32-byte transactions necessary to service all of the threads of the
warp} (BPG §10.2.1 \cite{nvidia2026cudabestpractices}); Lesson~4 counts
these sectors.}

\begin{note}
  Simulating a sectored cache changes what the simulator counts.  A
  block can be present with only some sectors valid, so an access can
  hit the tag and still miss the sector; a simulator that counts only
  tag hits overestimates the hit rate and underestimates DRAM traffic.
\end{note}

\subsection{Open-source GPU simulators}

\Cref{tab:l08-simulators} lists the open-source GPU simulators that
are widely used in research.  They differ in the ISA they model, in how
they obtain instructions, and in what they model beyond one GPU.

\begin{table}[tb]
  \centering
  \caption{Open-source GPU simulators, as classified in this chapter.}
  \label{tab:l08-simulators}
  \setlength{\tabcolsep}{4pt}
  \begin{tabular}{@{}>{\raggedright\arraybackslash}p{24mm}
                  >{\raggedright\arraybackslash}p{25mm}
                  >{\raggedright\arraybackslash}p{19mm}
                  >{\raggedright\arraybackslash}p{40mm}@{}}
    \toprule
    simulator & ISA & type & scope \\
    \midrule
    GPGPU-Sim \cite{bakhoda2009gpgpusim} & NVIDIA PTX, SASS
      & execution-driven & a GPU running general-purpose code \\
    Accel-Sim \cite{khairy2020accelsim} & NVIDIA PTX, SASS
      & trace-driven & GPUs and accelerators \\
    MGPUSim \cite{sun2019mgpusim} & AMD GPU & execution-driven
      & multiple GPUs \\
    MacSim \cite{gthparch2026macsim} & NVIDIA and Intel GPUs
      & trace-driven & heterogeneous CPU and GPU systems \\
    gem5-gpu \cite{power2015gem5gpu} & AMD GPU or NVIDIA PTX
      & execution-driven & heterogeneous CPU and GPU systems \\
    \bottomrule
  \end{tabular}
\end{table}

GPGPU-Sim is the most widely used of them.  gem5-gpu combines it with
gem5, the most widely used CPU simulator, so that a CPU and a GPU
execute together and share a memory system.  MGPUSim models AMD GPUs
and is the base on which the Photon sampling methodology of
\cref{sec:l08-accelerating} was built.\sidenote{GPGPU-Sim was introduced with a study of CUDA
workloads at ISPASS 2009 \cite{bakhoda2009gpgpusim}.  The table
follows the classification of this chapter's source.  In detail, Accel-Sim keeps GPGPU-Sim's
execution-driven PTX mode next to its trace-driven SASS mode
\cite{khairy2020accelsim}, and gem5-gpu, the
combination also listed as \enquote{Gem5-GPGPU-Sim}, couples gem5
with GPGPU-Sim's NVIDIA model; the AMD GPU model of gem5 is a separate
development.  MacSim's documentation describes it as configurable
as \enquote{either a trace driven or execution-driven cycle level
simulator} for x86, ARM64, PTX and Intel GEN instructions
\cite{gthparch2026macsim}.}

\begin{keypoint}
  A GPU simulator models warps with their active masks.  It either
  maintains the SIMT stack itself (execution-driven) or replays the
  recorded blocks and masks (trace-driven, tied to the trace machine's
  divergence handling).  Coalescing decides the number of transactions
  and is among the most important things to model; pre-coalesced
  traces fix the transaction size.  Sectored caches fill one sector per
  miss: less bandwidth, less implicit prefetching.
\end{keypoint}

\section{Analytical models of GPUs}
\label{sec:l08-analytical}

A cycle-level simulator answers a question about one design and one
program by running the program.\source{M8, analytical models; K12 ch.~2
Performance Principles, §4.1.1 Execution Time Modeling, §4.2.4
Roofline Model, §4.2.6 Other GPGPU Performance Modeling Techniques;
\cite{williams2009roofline,hong2009analytical,huang2014gpumech,karkhanis2004firstorder}.}
An analytical model answers it with a formula.

\begin{definition}[Analytical model]
  \label{def:l08-analytical-model}
  An \term{analytical model}[解析モデル] predicts performance from a
  small number of parameters of the design and the program, by
  equations rather than by executing the program.  A
  \term{first-order model}[一次近似モデル] keeps only the effects that
  dominate performance and ignores second-order interactions.
\end{definition}

Analytical models need no full execution, are simple enough to
reason about, and give insight: they say \emph{why} a design is fast or
slow and which parameter to change.  Their accuracy is lower than that
of a simulator, so they are used to explore a design space before
simulating its best points, and to explain simulation results.

\subsection{A first-order GPU design}

We size a GPU for one target: \num{e12} dot-product steps per second,
each step being \lstinline[language=C++]!sum += x[i] * y[i]!.

\begin{definition}[Fused multiply-add]
  \label{def:l08-fma}
  A \term{fused multiply-add (FMA)}[融合積和演算（FMA）]\index{FMA|see{fused multiply-add (FMA)}} computes
  $a \cdot b + c$ as one instruction with a single rounding.  It is
  counted as two floating-point operations (FLOPs).
\end{definition}

\begin{description}
  \item[Compute.]  Each step is one multiply and one add, so the target
    is \SI{2}{\tera\flop\per\second}, or \num{e12} FMAs per second.  At
    \SI{1}{\giga\hertz} this needs \num{1000} FMA units, at
    \SI{2}{\giga\hertz} \num{500}.
  \item[Memory.]  Each step reads one element of \texttt{x} and one of
    \texttt{y}.  Counting one byte per element, as this design
    exercise does, the GPU needs \SI{2}{\tera\byte\per\second} of
    memory bandwidth.\caution{Lecture vs reality: the
    \enquote{two bytes per step} assume one-byte elements.  With
    FP32 elements a step reads \SI{8}{\byte} for 2 FLOPs, and the same
    design needs \SI{8}{\tera\byte\per\second}, about 2.7 times the
    \SI{3}{\tera\byte\per\second} of an H100 SXM5 (preliminary
    figure, \cite{nvidia2022hopper}).}
  \item[Warps.]  \num{500} FMA units correspond to about 16 warps of
    width 32 ($16 \times 32 = 512$) issuing one FMA per lane per
    cycle.  The GPU can be built as 16~SMs that each issue one warp
    per cycle, or as 8~SMs that each issue two, both at
    \SI{2}{\giga\hertz}.
\end{description}

\subsection{Width and depth of multithreading}

How many warps an SM must hold follows from two parameters.

\begin{definition}[$W_{\text{width}}$ and $W_{\text{depth}}$]
  \label{def:l08-wwidth-wdepth}
  $W_{\text{width}}$ is the number of threads that execute in one
  cycle (for example 32, one warp).  $W_{\text{depth}}$ is the number of
  warps that the SM can schedule while one warp stalls, so that the
  stall is hidden.  An SM holds
  $W_{\text{width}} \times W_{\text{depth}}$ resident threads.
\end{definition}

\begin{figure}[tb]
  \centering
  % W_width and W_depth (Lesson 8).  Left: issue slots of one SM for
% 1, 2 and 4 warps, each warp issuing one compute (C) and one memory (M)
% instruction and then waiting until 8 cycles after its compute
% instruction; grey = idle issue slot.  With 4 warps the stall is fully
% hidden (W_depth = 4).  Right: W_width x W_depth resident threads.
% Redraws the outline figure m8l3_arch2 (our own latency of 8 cycles).
% Usage: % W_width and W_depth (Lesson 8).  Left: issue slots of one SM for
% 1, 2 and 4 warps, each warp issuing one compute (C) and one memory (M)
% instruction and then waiting until 8 cycles after its compute
% instruction; grey = idle issue slot.  With 4 warps the stall is fully
% hidden (W_depth = 4).  Right: W_width x W_depth resident threads.
% Redraws the outline figure m8l3_arch2 (our own latency of 8 cycles).
% Usage: % W_width and W_depth (Lesson 8).  Left: issue slots of one SM for
% 1, 2 and 4 warps, each warp issuing one compute (C) and one memory (M)
% instruction and then waiting until 8 cycles after its compute
% instruction; grey = idle issue slot.  With 4 warps the stall is fully
% hidden (W_depth = 4).  Right: W_width x W_depth resident threads.
% Redraws the outline figure m8l3_arch2 (our own latency of 8 cycles).
% Usage: \input{figures/l08-wdepth}
% Width: body (120 mm); the drawing is about 116 mm wide.
\begin{tikzpicture}[x=4.4mm, y=4.4mm, font=\footnotesize\sffamily]
  % \slots{y}{label}{list}: a row of issue slots; "-" is idle
  \def\idle{-}%
  \newcommand{\slots}[3]{%
    \node[anchor=east, inner sep=1pt] at (-0.2,#1-0.5) {#2};
    \foreach \s [count=\i from 0] in {#3} {
      \ifx\s\idle
        \draw[fill=ln-rule!40, draw=ln-rule] (\i,#1) rectangle ++(1,-1);
      \else
        \draw[fill=ln-box] (\i,#1) rectangle ++(1,-1) node[midway] {\s};
      \fi
    }}
  \slots{0}{\iflangja{1ワープ}{1 warp}}{C1,M1,-,-,-,-,-,-,C1,M1,-,-,-,-,-,-}
  \slots{-1.6}{\iflangja{2ワープ}{2 warps}}{C1,M1,C2,M2,-,-,-,-,C1,M1,C2,M2,-,-,-,-}
  \slots{-3.2}{\iflangja{4ワープ}{4 warps}}{C1,M1,C2,M2,C3,M3,C4,M4,C1,M1,C2,M2,C3,M3,C4,M4}
  \draw[->, ln-muted] (0,-4.8) -- (16.3,-4.8)
    node[below left, inner sep=1pt] {\iflangja{サイクル}{cycle}};
  \draw[ln-caution, thick, ->] (0,0.4) -- (8,0.4)
    node[midway, above, inner sep=1pt, text=ln-caution]
    {\iflangja{ワープ1の周期 = 8サイクル}{one round of warp 1 = 8 cycles}};
  % W_width x W_depth grid
  \begin{scope}[shift={(18.6,-4.2)}, x=1.4mm, y=1.4mm]
    \foreach \r in {0,...,3} {
      \foreach \c in {0,...,7} {
        \draw[fill=ln-accent!15] (\c,\r) circle (0.45);
      }
    }
    \draw[->] (0,4.3) -- (7.6,4.3) node[midway, above, inner sep=1pt]
      {$W_{\text{width}}$};
    \draw[->] (8.4,3.4) -- (8.4,-0.4) node[midway, right, inner sep=1pt]
      {$W_{\text{depth}}$};
    \node[anchor=north, align=center, text=ln-muted] at (4,-0.8)
      {\iflangja{常駐スレッド数\\$= W_{\text{width}} \times W_{\text{depth}}$}%
                {resident threads\\$= W_{\text{width}} \times W_{\text{depth}}$}};
  \end{scope}
\end{tikzpicture}

% Width: body (120 mm); the drawing is about 116 mm wide.
\begin{tikzpicture}[x=4.4mm, y=4.4mm, font=\footnotesize\sffamily]
  % \slots{y}{label}{list}: a row of issue slots; "-" is idle
  \def\idle{-}%
  \newcommand{\slots}[3]{%
    \node[anchor=east, inner sep=1pt] at (-0.2,#1-0.5) {#2};
    \foreach \s [count=\i from 0] in {#3} {
      \ifx\s\idle
        \draw[fill=ln-rule!40, draw=ln-rule] (\i,#1) rectangle ++(1,-1);
      \else
        \draw[fill=ln-box] (\i,#1) rectangle ++(1,-1) node[midway] {\s};
      \fi
    }}
  \slots{0}{\iflangja{1ワープ}{1 warp}}{C1,M1,-,-,-,-,-,-,C1,M1,-,-,-,-,-,-}
  \slots{-1.6}{\iflangja{2ワープ}{2 warps}}{C1,M1,C2,M2,-,-,-,-,C1,M1,C2,M2,-,-,-,-}
  \slots{-3.2}{\iflangja{4ワープ}{4 warps}}{C1,M1,C2,M2,C3,M3,C4,M4,C1,M1,C2,M2,C3,M3,C4,M4}
  \draw[->, ln-muted] (0,-4.8) -- (16.3,-4.8)
    node[below left, inner sep=1pt] {\iflangja{サイクル}{cycle}};
  \draw[ln-caution, thick, ->] (0,0.4) -- (8,0.4)
    node[midway, above, inner sep=1pt, text=ln-caution]
    {\iflangja{ワープ1の周期 = 8サイクル}{one round of warp 1 = 8 cycles}};
  % W_width x W_depth grid
  \begin{scope}[shift={(18.6,-4.2)}, x=1.4mm, y=1.4mm]
    \foreach \r in {0,...,3} {
      \foreach \c in {0,...,7} {
        \draw[fill=ln-accent!15] (\c,\r) circle (0.45);
      }
    }
    \draw[->] (0,4.3) -- (7.6,4.3) node[midway, above, inner sep=1pt]
      {$W_{\text{width}}$};
    \draw[->] (8.4,3.4) -- (8.4,-0.4) node[midway, right, inner sep=1pt]
      {$W_{\text{depth}}$};
    \node[anchor=north, align=center, text=ln-muted] at (4,-0.8)
      {\iflangja{常駐スレッド数\\$= W_{\text{width}} \times W_{\text{depth}}$}%
                {resident threads\\$= W_{\text{width}} \times W_{\text{depth}}$}};
  \end{scope}
\end{tikzpicture}

% Width: body (120 mm); the drawing is about 116 mm wide.
\begin{tikzpicture}[x=4.4mm, y=4.4mm, font=\footnotesize\sffamily]
  % \slots{y}{label}{list}: a row of issue slots; "-" is idle
  \def\idle{-}%
  \newcommand{\slots}[3]{%
    \node[anchor=east, inner sep=1pt] at (-0.2,#1-0.5) {#2};
    \foreach \s [count=\i from 0] in {#3} {
      \ifx\s\idle
        \draw[fill=ln-rule!40, draw=ln-rule] (\i,#1) rectangle ++(1,-1);
      \else
        \draw[fill=ln-box] (\i,#1) rectangle ++(1,-1) node[midway] {\s};
      \fi
    }}
  \slots{0}{\iflangja{1ワープ}{1 warp}}{C1,M1,-,-,-,-,-,-,C1,M1,-,-,-,-,-,-}
  \slots{-1.6}{\iflangja{2ワープ}{2 warps}}{C1,M1,C2,M2,-,-,-,-,C1,M1,C2,M2,-,-,-,-}
  \slots{-3.2}{\iflangja{4ワープ}{4 warps}}{C1,M1,C2,M2,C3,M3,C4,M4,C1,M1,C2,M2,C3,M3,C4,M4}
  \draw[->, ln-muted] (0,-4.8) -- (16.3,-4.8)
    node[below left, inner sep=1pt] {\iflangja{サイクル}{cycle}};
  \draw[ln-caution, thick, ->] (0,0.4) -- (8,0.4)
    node[midway, above, inner sep=1pt, text=ln-caution]
    {\iflangja{ワープ1の周期 = 8サイクル}{one round of warp 1 = 8 cycles}};
  % W_width x W_depth grid
  \begin{scope}[shift={(18.6,-4.2)}, x=1.4mm, y=1.4mm]
    \foreach \r in {0,...,3} {
      \foreach \c in {0,...,7} {
        \draw[fill=ln-accent!15] (\c,\r) circle (0.45);
      }
    }
    \draw[->] (0,4.3) -- (7.6,4.3) node[midway, above, inner sep=1pt]
      {$W_{\text{width}}$};
    \draw[->] (8.4,3.4) -- (8.4,-0.4) node[midway, right, inner sep=1pt]
      {$W_{\text{depth}}$};
    \node[anchor=north, align=center, text=ln-muted] at (4,-0.8)
      {\iflangja{常駐スレッド数\\$= W_{\text{width}} \times W_{\text{depth}}$}%
                {resident threads\\$= W_{\text{width}} \times W_{\text{depth}}$}};
  \end{scope}
\end{tikzpicture}

  \caption{Hiding a memory stall.  Each warp issues a compute (C) and a
    memory (M) instruction and can issue again only 8~cycles after its
    compute instruction.  One warp keeps the issue slot busy 2 cycles
    in 8, two warps 4 in 8, four warps all the time:
    $W_{\text{depth}} = 4$.}
  \label{fig:l08-wdepth}
\end{figure}

\Cref{fig:l08-wdepth} shows the idea for a program whose warps
alternate one compute and one memory instruction.  With one warp, the
SM issues two instructions and then idles until the memory instruction
completes; with two warps, the second warp fills two of the idle
slots; with four, the idle slots disappear.  In general the stall must
be covered by the work of the other warps: if each warp issues $n$
instructions and then cannot issue again until $L$ cycles after the
first of them, about
\begin{equation}
  W_{\text{depth}} \approx \frac{L}{n}
  \label{eq:l08-wdepth}
\end{equation}
warps are needed.\tip{Count the cycles a warp keeps the issue slot busy
between two waits, and divide the wait by it.}  With a memory latency
of \num{200}~cycles and one compute plus one memory instruction per
warp, $W_{\text{depth}} = 200/(1+1) = 100$; with one memory instruction
per four compute instructions, $200/(1+4) = 40$.  The more compute per
memory access, the fewer warps are needed to hide the latency.

Both parameters are set by hardware.  $W_{\text{width}}$ is bounded by
the number of ALUs (and the width of the scheduler).
$W_{\text{depth}}$ is bounded by the storage for the warps' state: a
$W_{\text{depth}}$ of 20 needs $20 \times 32 \times r$ registers for $r$
registers per thread, and at least 20 program counters, one per
warp.\margin{Real limits: at most 64 resident warps per SM
(compute capability 7.0, 8.0, 9.0) and \num{65536} 32-bit registers per
SM, so 64 warps of 32 threads get at most 32 registers each (CPG
Table~21 \cite{nvidia2024cudaguide}).}  Equation~\eqref{eq:l08-wdepth}
with a latency of hundreds of cycles explains why GPUs keep dozens of
warps per SM and why their register files are so large (Lesson~6).

\subsection{How many SMs, and sub-cores}

The two designs above, 16~SMs issuing one warp per cycle and 8~SMs
issuing two, have the same execution bandwidth and, for the same chip
area, the same total cache capacity.  They differ in how that capacity
is split.  Fewer, larger SMs have larger caches with longer access
times, but the thread blocks on one SM share more data, which can raise
the hit rate; more, smaller SMs have smaller, faster caches.  Which is
better depends on how much data the blocks share, a trade-off that a
first-order model cannot settle; it needs a finer model or a
simulation.

\begin{definition}[Sub-core]
  \label{def:l08-sub-core}
  A \term{sub-core}[サブコア] is one of several partitions of an SM,
  each with its own program counters, decoder, warp scheduler and
  execution lanes, which share the SM's instruction cache and, through
  a crossbar, its large caches and shared memory
  (\cref{fig:l08-subcore}).
\end{definition}

\begin{figure}[tb]
  \centering
  % An SM divided into sub-cores (Lesson 8).  Each sub-core has its own
% PC registers, decoder, warp scheduler and execution lanes; all
% sub-cores share the instruction cache at the top and, through a
% crossbar, the large caches at the bottom.
% Redraws the outline figure m8l3_subcore (four sub-cores instead of
% three, as in current NVIDIA SMs).
% Usage: % An SM divided into sub-cores (Lesson 8).  Each sub-core has its own
% PC registers, decoder, warp scheduler and execution lanes; all
% sub-cores share the instruction cache at the top and, through a
% crossbar, the large caches at the bottom.
% Redraws the outline figure m8l3_subcore (four sub-cores instead of
% three, as in current NVIDIA SMs).
% Usage: % An SM divided into sub-cores (Lesson 8).  Each sub-core has its own
% PC registers, decoder, warp scheduler and execution lanes; all
% sub-cores share the instruction cache at the top and, through a
% crossbar, the large caches at the bottom.
% Redraws the outline figure m8l3_subcore (four sub-cores instead of
% three, as in current NVIDIA SMs).
% Usage: \input{figures/l08-subcore}
% Width: body (120 mm); the drawing is about 108 mm wide.
\begin{tikzpicture}[x=1mm, y=1mm, font=\scriptsize\sffamily,
    u/.style={draw, fill=white, minimum width=21mm, minimum height=4.6mm,
              inner sep=0.5pt}]
  \draw[ln-accent, thick, rounded corners=1.5mm] (0,-2) rectangle (108,50);
  \node[anchor=north west, text=ln-accent, font=\scriptsize\sffamily\bfseries]
    at (1,49) {SM};
  \node[draw, fill=ln-tip!15, minimum width=96mm, minimum height=5mm]
    at (57,43) {\iflangja{命令キャッシュ（共有）}{instruction cache (shared)}};
  \foreach \k in {0,...,3} {
    \pgfmathsetmacro{\x}{8+\k*25}
    \draw[fill=ln-accent!10] (\x,11) rectangle ++(23,28);
    \node[anchor=north] at (\x+11.5,38.6) {\iflangja{サブコア}{sub-core} \k};
    \node[u] at (\x+11.5,31) {\iflangja{PC レジスタ}{PC registers}};
    \node[u] at (\x+11.5,26) {\iflangja{デコーダ}{decoder}};
    \node[u] at (\x+11.5,21) {\iflangja{ワープスケジューラ}{warp scheduler}};
    \node[u, fill=ln-caution!12] at (\x+11.5,15.6) {\iflangja{実行レーン}{lanes (SPs)}};
  }
  \node[draw, fill=ln-box, minimum width=96mm, minimum height=4.5mm]
    at (57,6) {\iflangja{クロスバ}{crossbar}};
  \node[draw, fill=ln-tip!15, minimum width=96mm, minimum height=4.5mm]
    at (57,0.8) {\iflangja{共有メモリと L1 キャッシュ（共有）}{shared memory and L1 cache (shared)}};
\end{tikzpicture}

% Width: body (120 mm); the drawing is about 108 mm wide.
\begin{tikzpicture}[x=1mm, y=1mm, font=\scriptsize\sffamily,
    u/.style={draw, fill=white, minimum width=21mm, minimum height=4.6mm,
              inner sep=0.5pt}]
  \draw[ln-accent, thick, rounded corners=1.5mm] (0,-2) rectangle (108,50);
  \node[anchor=north west, text=ln-accent, font=\scriptsize\sffamily\bfseries]
    at (1,49) {SM};
  \node[draw, fill=ln-tip!15, minimum width=96mm, minimum height=5mm]
    at (57,43) {\iflangja{命令キャッシュ（共有）}{instruction cache (shared)}};
  \foreach \k in {0,...,3} {
    \pgfmathsetmacro{\x}{8+\k*25}
    \draw[fill=ln-accent!10] (\x,11) rectangle ++(23,28);
    \node[anchor=north] at (\x+11.5,38.6) {\iflangja{サブコア}{sub-core} \k};
    \node[u] at (\x+11.5,31) {\iflangja{PC レジスタ}{PC registers}};
    \node[u] at (\x+11.5,26) {\iflangja{デコーダ}{decoder}};
    \node[u] at (\x+11.5,21) {\iflangja{ワープスケジューラ}{warp scheduler}};
    \node[u, fill=ln-caution!12] at (\x+11.5,15.6) {\iflangja{実行レーン}{lanes (SPs)}};
  }
  \node[draw, fill=ln-box, minimum width=96mm, minimum height=4.5mm]
    at (57,6) {\iflangja{クロスバ}{crossbar}};
  \node[draw, fill=ln-tip!15, minimum width=96mm, minimum height=4.5mm]
    at (57,0.8) {\iflangja{共有メモリと L1 キャッシュ（共有）}{shared memory and L1 cache (shared)}};
\end{tikzpicture}

% Width: body (120 mm); the drawing is about 108 mm wide.
\begin{tikzpicture}[x=1mm, y=1mm, font=\scriptsize\sffamily,
    u/.style={draw, fill=white, minimum width=21mm, minimum height=4.6mm,
              inner sep=0.5pt}]
  \draw[ln-accent, thick, rounded corners=1.5mm] (0,-2) rectangle (108,50);
  \node[anchor=north west, text=ln-accent, font=\scriptsize\sffamily\bfseries]
    at (1,49) {SM};
  \node[draw, fill=ln-tip!15, minimum width=96mm, minimum height=5mm]
    at (57,43) {\iflangja{命令キャッシュ（共有）}{instruction cache (shared)}};
  \foreach \k in {0,...,3} {
    \pgfmathsetmacro{\x}{8+\k*25}
    \draw[fill=ln-accent!10] (\x,11) rectangle ++(23,28);
    \node[anchor=north] at (\x+11.5,38.6) {\iflangja{サブコア}{sub-core} \k};
    \node[u] at (\x+11.5,31) {\iflangja{PC レジスタ}{PC registers}};
    \node[u] at (\x+11.5,26) {\iflangja{デコーダ}{decoder}};
    \node[u] at (\x+11.5,21) {\iflangja{ワープスケジューラ}{warp scheduler}};
    \node[u, fill=ln-caution!12] at (\x+11.5,15.6) {\iflangja{実行レーン}{lanes (SPs)}};
  }
  \node[draw, fill=ln-box, minimum width=96mm, minimum height=4.5mm]
    at (57,6) {\iflangja{クロスバ}{crossbar}};
  \node[draw, fill=ln-tip!15, minimum width=96mm, minimum height=4.5mm]
    at (57,0.8) {\iflangja{共有メモリと L1 キャッシュ（共有）}{shared memory and L1 cache (shared)}};
\end{tikzpicture}

  \caption{An SM divided into four sub-cores.  Each has its own front
    end and lanes; the instruction cache and the data caches are
    shared.}
  \label{fig:l08-subcore}
\end{figure}

Sub-cores are a middle way between the two designs: each schedules a
small number of warps over its own lanes, while large structures stay
shared.  For a model this adds resource management: two warps on
different sub-cores compete for the shared caches and crossbar but not
for an issue slot.\margin{Current NVIDIA SMs have 128 FP32 lanes
split into four partitions \cite{nvidia2022hopper}, so each sub-core
issues for 32 lanes.}

\begin{keypoint}
  A first-order design starts from the target throughput: FLOPs fix the
  number of FMA units (at a given clock), bytes per FLOP fix the
  memory bandwidth, and the latency to hide fixes $W_{\text{depth}}
  \approx L/n$, hence the registers and program counters per SM.
\end{keypoint}

\subsection{The Roofline model}

The first-order design asked how much compute and how much bandwidth a
program needs.  The Roofline model asks the reverse question: given a
machine, which of the two limits a program.

\begin{definition}[Arithmetic intensity]
  \label{def:l08-arithmetic-intensity}
  The \term{arithmetic intensity}[演算強度] of a program or kernel is
  the number of floating-point operations it performs per byte of
  memory it accesses, in FLOP/B.
\end{definition}

\begin{definition}[Roofline model]
  \label{def:l08-roofline}
  The \term{Roofline model}[ルーフラインモデル] bounds the attainable
  performance of a kernel with arithmetic intensity $I$ on a machine
  with peak compute performance $P_{\text{peak}}$ and peak memory
  bandwidth $B_{\text{peak}}$ by
  \begin{equation}
    P(I) = \min\bigl(P_{\text{peak}},\; B_{\text{peak}} \cdot I\bigr).
    \label{eq:l08-roofline}
  \end{equation}
  On log-log axes the bound is a slanted line (slope one) followed by a
  flat roof (\cref{fig:l08-roofline}).  A kernel left of the
  \term{ridge point}[リッジポイント], the intensity
  $I^{*} = P_{\text{peak}}/B_{\text{peak}}$ at which the two lines meet,
  is \term{memory-bound}[メモリ律速]: its performance is limited by
  bandwidth.  A kernel right of it is
  \term{compute-bound}[演算律速]: limited by the peak FLOP rate.
\end{definition}

\begin{figure}[tb]
  \centering
  % The Roofline model (Lesson 8), log-log axes, for our example machine
% with a peak of 16 TFLOP/s and 1 TB/s of memory bandwidth (ridge point
% 16 FLOP/B).  The slanted roof is bandwidth x intensity (memory-bound
% region), the flat roof the peak (compute-bound region).  Kernel A has
% an arithmetic intensity of 0.5 FLOP/B, kernel B of 2 FLOP/B: raising
% the intensity moves a kernel right, toward the ridge point.
% Redraws the outline figure m8l3_roffline [sic] with our own numbers.
% Usage: % The Roofline model (Lesson 8), log-log axes, for our example machine
% with a peak of 16 TFLOP/s and 1 TB/s of memory bandwidth (ridge point
% 16 FLOP/B).  The slanted roof is bandwidth x intensity (memory-bound
% region), the flat roof the peak (compute-bound region).  Kernel A has
% an arithmetic intensity of 0.5 FLOP/B, kernel B of 2 FLOP/B: raising
% the intensity moves a kernel right, toward the ridge point.
% Redraws the outline figure m8l3_roffline [sic] with our own numbers.
% Usage: % The Roofline model (Lesson 8), log-log axes, for our example machine
% with a peak of 16 TFLOP/s and 1 TB/s of memory bandwidth (ridge point
% 16 FLOP/B).  The slanted roof is bandwidth x intensity (memory-bound
% region), the flat roof the peak (compute-bound region).  Kernel A has
% an arithmetic intensity of 0.5 FLOP/B, kernel B of 2 FLOP/B: raising
% the intensity moves a kernel right, toward the ridge point.
% Redraws the outline figure m8l3_roffline [sic] with our own numbers.
% Usage: \input{figures/l08-roofline}
% Width: body (120 mm); the drawing is about 110 mm wide.
\begin{tikzpicture}[x=8.5mm, y=6mm, font=\footnotesize\sffamily,
    >={Stealth[length=1.5mm]}]
  % regions (x = log2 FLOP/B from -3 to 7, y = log2 TFLOP/s from -3 to 5)
  \fill[ln-accent!8] (-3,-3) rectangle (4,5.4);
  \fill[ln-tip!10] (4,-3) rectangle (7,5.4);
  \node[text=ln-accent, align=center] at (-0.2,4.6)
    {\iflangja{メモリ律速}{memory-bound}};
  \node[text=ln-tip, align=center] at (5.5,0)
    {\iflangja{演算律速}{compute-bound}};
  % axes
  \draw[->] (-3,-3) -- (7.3,-3);
  \draw[->] (-3,-3) -- (-3,5.6);
  \foreach \e/\l in {-3/{1/8},-2/{1/4},-1/{1/2},0/1,1/2,2/4,3/8,4/16,5/32,6/64,7/128} {
    \draw (\e,-3) -- ++(0,-0.15) node[below, inner sep=1pt] {\l};
  }
  \foreach \e/\l in {-3/{1/8},-1/{1/2},1/2,3/8,4/16,5/32} {
    \draw (-3,\e) -- ++(-0.12,0) node[left, inner sep=1pt] {\l};
  }
  \node[below, inner sep=1pt] at (2,-4.1)
    {\iflangja{演算強度（FLOP/B，対数目盛）}{arithmetic intensity (FLOP/B, log scale)}};
  \node[rotate=90, above, inner sep=1pt] at (-4.15,1)
    {\iflangja{性能（TFLOP/s，対数目盛）}{performance (TFLOP/s, log scale)}};
  % roofs
  \draw[very thick, ln-accent] (-3,-3) -- (4,4);
  \draw[very thick, ln-tip] (4,4) -- (7,4);
  \node[above, text=ln-tip] at (5.5,4)
    {\iflangja{ピーク 16\,TFLOP/s}{peak 16\,TFLOP/s}};
  \node[rotate=35.5, below, text=ln-accent] at (2.7,2.7)
    {\iflangja{帯域 1\,TB/s $\times$ 演算強度}{1\,TB/s $\times$ intensity}};
  % ridge point
  \draw[dashed, ln-muted] (4,-3) -- (4,4);
  \fill[ln-caution] (4,4) circle (0.09);
  \node[right, text=ln-caution, align=left] at (4.1,3.25)
    {\iflangja{リッジポイント\\16\,FLOP/B}{ridge point\\16\,FLOP/B}};
  % kernels
  \fill (-1,-1) circle (0.1) node[below right, inner sep=1.5pt] {A: 0.5};
  \fill (1,1) circle (0.1) node[below right, inner sep=1.5pt] {B: 2};
  \draw[->, thick, ln-caution] (-0.85,-0.75) to[bend left=20] (0.8,0.85);
\end{tikzpicture}

% Width: body (120 mm); the drawing is about 110 mm wide.
\begin{tikzpicture}[x=8.5mm, y=6mm, font=\footnotesize\sffamily,
    >={Stealth[length=1.5mm]}]
  % regions (x = log2 FLOP/B from -3 to 7, y = log2 TFLOP/s from -3 to 5)
  \fill[ln-accent!8] (-3,-3) rectangle (4,5.4);
  \fill[ln-tip!10] (4,-3) rectangle (7,5.4);
  \node[text=ln-accent, align=center] at (-0.2,4.6)
    {\iflangja{メモリ律速}{memory-bound}};
  \node[text=ln-tip, align=center] at (5.5,0)
    {\iflangja{演算律速}{compute-bound}};
  % axes
  \draw[->] (-3,-3) -- (7.3,-3);
  \draw[->] (-3,-3) -- (-3,5.6);
  \foreach \e/\l in {-3/{1/8},-2/{1/4},-1/{1/2},0/1,1/2,2/4,3/8,4/16,5/32,6/64,7/128} {
    \draw (\e,-3) -- ++(0,-0.15) node[below, inner sep=1pt] {\l};
  }
  \foreach \e/\l in {-3/{1/8},-1/{1/2},1/2,3/8,4/16,5/32} {
    \draw (-3,\e) -- ++(-0.12,0) node[left, inner sep=1pt] {\l};
  }
  \node[below, inner sep=1pt] at (2,-4.1)
    {\iflangja{演算強度（FLOP/B，対数目盛）}{arithmetic intensity (FLOP/B, log scale)}};
  \node[rotate=90, above, inner sep=1pt] at (-4.15,1)
    {\iflangja{性能（TFLOP/s，対数目盛）}{performance (TFLOP/s, log scale)}};
  % roofs
  \draw[very thick, ln-accent] (-3,-3) -- (4,4);
  \draw[very thick, ln-tip] (4,4) -- (7,4);
  \node[above, text=ln-tip] at (5.5,4)
    {\iflangja{ピーク 16\,TFLOP/s}{peak 16\,TFLOP/s}};
  \node[rotate=35.5, below, text=ln-accent] at (2.7,2.7)
    {\iflangja{帯域 1\,TB/s $\times$ 演算強度}{1\,TB/s $\times$ intensity}};
  % ridge point
  \draw[dashed, ln-muted] (4,-3) -- (4,4);
  \fill[ln-caution] (4,4) circle (0.09);
  \node[right, text=ln-caution, align=left] at (4.1,3.25)
    {\iflangja{リッジポイント\\16\,FLOP/B}{ridge point\\16\,FLOP/B}};
  % kernels
  \fill (-1,-1) circle (0.1) node[below right, inner sep=1.5pt] {A: 0.5};
  \fill (1,1) circle (0.1) node[below right, inner sep=1.5pt] {B: 2};
  \draw[->, thick, ln-caution] (-0.85,-0.75) to[bend left=20] (0.8,0.85);
\end{tikzpicture}

% Width: body (120 mm); the drawing is about 110 mm wide.
\begin{tikzpicture}[x=8.5mm, y=6mm, font=\footnotesize\sffamily,
    >={Stealth[length=1.5mm]}]
  % regions (x = log2 FLOP/B from -3 to 7, y = log2 TFLOP/s from -3 to 5)
  \fill[ln-accent!8] (-3,-3) rectangle (4,5.4);
  \fill[ln-tip!10] (4,-3) rectangle (7,5.4);
  \node[text=ln-accent, align=center] at (-0.2,4.6)
    {\iflangja{メモリ律速}{memory-bound}};
  \node[text=ln-tip, align=center] at (5.5,0)
    {\iflangja{演算律速}{compute-bound}};
  % axes
  \draw[->] (-3,-3) -- (7.3,-3);
  \draw[->] (-3,-3) -- (-3,5.6);
  \foreach \e/\l in {-3/{1/8},-2/{1/4},-1/{1/2},0/1,1/2,2/4,3/8,4/16,5/32,6/64,7/128} {
    \draw (\e,-3) -- ++(0,-0.15) node[below, inner sep=1pt] {\l};
  }
  \foreach \e/\l in {-3/{1/8},-1/{1/2},1/2,3/8,4/16,5/32} {
    \draw (-3,\e) -- ++(-0.12,0) node[left, inner sep=1pt] {\l};
  }
  \node[below, inner sep=1pt] at (2,-4.1)
    {\iflangja{演算強度（FLOP/B，対数目盛）}{arithmetic intensity (FLOP/B, log scale)}};
  \node[rotate=90, above, inner sep=1pt] at (-4.15,1)
    {\iflangja{性能（TFLOP/s，対数目盛）}{performance (TFLOP/s, log scale)}};
  % roofs
  \draw[very thick, ln-accent] (-3,-3) -- (4,4);
  \draw[very thick, ln-tip] (4,4) -- (7,4);
  \node[above, text=ln-tip] at (5.5,4)
    {\iflangja{ピーク 16\,TFLOP/s}{peak 16\,TFLOP/s}};
  \node[rotate=35.5, below, text=ln-accent] at (2.7,2.7)
    {\iflangja{帯域 1\,TB/s $\times$ 演算強度}{1\,TB/s $\times$ intensity}};
  % ridge point
  \draw[dashed, ln-muted] (4,-3) -- (4,4);
  \fill[ln-caution] (4,4) circle (0.09);
  \node[right, text=ln-caution, align=left] at (4.1,3.25)
    {\iflangja{リッジポイント\\16\,FLOP/B}{ridge point\\16\,FLOP/B}};
  % kernels
  \fill (-1,-1) circle (0.1) node[below right, inner sep=1.5pt] {A: 0.5};
  \fill (1,1) circle (0.1) node[below right, inner sep=1.5pt] {B: 2};
  \draw[->, thick, ln-caution] (-0.85,-0.75) to[bend left=20] (0.8,0.85);
\end{tikzpicture}

  \caption{The Roofline model of a machine with
    \SI{16}{\tera\flop\per\second} and \SI{1}{\tera\byte\per\second}
    (ridge point \SI{16}{\flop\per\byte}).  Kernel A
    (\SI{0.5}{\flop\per\byte}) reaches at most \SI{0.5}{\tera\flop\per\second};
    raising the intensity to \SI{2}{\flop\per\byte} (kernel B) quadruples the
    bound.}
  \label{fig:l08-roofline}
\end{figure}

Williams et al., who introduced the model, measure the intensity
after the caches: their
\term{operational intensity}[演算強度（DRAMバイト当たり）] counts the
operations per byte of DRAM traffic, so a kernel whose data stay in the
caches has a higher operational than arithmetic intensity
\cite{williams2009roofline}.\caution{Lecture vs paper: this chapter,
like the lecture, counts the bytes the program requests.  The paper
uses \enquote{the term operational intensity to mean operations per byte
of DRAM traffic} and states the bound as \enquote{Attainable GFlops/sec
= Min(Peak Floating Point Performance, Peak Memory Bandwidth x
Operational Intensity)} \cite{williams2009roofline}.}

Two examples in the units of the first-order design, one byte per
element, show how a kernel moves on the chart.  A vector sum
\mbox{\lstinline[language=C++]!z[i] = x[i] + y[i]!}, counting only its two
loads (one byte each), performs 1~FLOP per 2~bytes: its intensity is
\SI{0.5}{\flop\per\byte}, and the kernel is memory-bound.  The variant
\mbox{\lstinline[language=C++]!sum += x[i]*x[i]*y[i]*y[i]!} of the
dot product reads the same
2~bytes but performs 4~FLOPs (three multiplications and one addition),
an intensity of \SI{2}{\flop\per\byte}: it moves right on the chart,
toward the compute-bound region.  The model's
practical use is to judge whether a kernel has room to improve: a
kernel far below its roof is limited by something else (latency,
divergence, bank conflicts), one on the roof can only be improved by
raising its intensity or by moving to another machine.  The roofs are
drawn once per machine, \enquote{not once per kernel}
\cite{williams2009roofline}; each kernel is then a point on the
intensity axis.\margin{Ridge points of real GPUs (FP16 tensor peak divided by
HBM bandwidth): V100 \SI{138.9}{\flop\per\byte} \cite{nvidia2026gemmguide};
A100 $312/1.555 \approx 201$; H100 SXM5
$1000/3 \approx 333$ (preliminary figures \cite{nvidia2020ampere,nvidia2022hopper}).}

\begin{example}
  \label{ex:l08-roofline}
  Our machine of \cref{fig:l08-roofline} has
  $P_{\text{peak}} = \SI{16}{\tera\flop\per\second}$ and
  $B_{\text{peak}} = \SI{1}{\tera\byte\per\second}$, so
  $I^{*} = 16/1 = \SI{16}{\flop\per\byte}$.  An FP32 dot product reads
  \SI{8}{\byte} per FMA (2~FLOPs), $I = 2/8 = \SI{0.25}{\flop\per\byte}$, and
  attains at most $1 \times 0.25 = \SI{0.25}{\tera\flop\per\second}$,
  about 1.6\,\% of the peak; the four-FLOP variant above reaches
  \SI{0.5}{\flop\per\byte} in FP32.  A kernel with $I = \SI{32}{\flop\per\byte}$, such
  as a large tiled matrix multiplication, attains
  $\min(16, 32) = \SI{16}{\tera\flop\per\second}$, the full
  peak.
\end{example}

\subsection{CPI models}

A second family of analytical models starts from the time per
instruction.

\begin{definition}[CPI]
  \label{def:l08-cpi}
  The \term{cycles per instruction (CPI)}[命令当たりサイクル数（CPI）]\index{CPI|see{cycles per instruction (CPI)}}
  of a program on a processor is the number of cycles it takes divided
  by the number of instructions it executes, the reciprocal of the IPC
  of Lesson~1.
\end{definition}

The simplest CPI model adds to the CPI of an ideal pipeline the cost of
every kind of disruptive event, weighted by how often it occurs:
\begin{equation}
  \mathrm{CPI} = \mathrm{CPI}_{\text{steady}} +
    \sum_{i} f_i \cdot p_i ,
  \label{eq:l08-cpi}
\end{equation}
where $f_i$ is the number of events of kind $i$ per instruction and
$p_i$ the penalty of one such event in cycles.  The model assumes that
penalties do not overlap: they are serialised.\caution{The serialised
sum is pessimistic for out-of-order and multithreaded processors, which
overlap penalties with each other and with useful work.}

\begin{example}
  \label{ex:l08-cpi}
  A five-stage in-order pipeline has $\mathrm{CPI}_{\text{steady}} = 1$.
  Branch mispredictions cost \num{3}~cycles and occur in 2\,\% of the
  instructions, cache misses cost \num{5}~cycles and occur in 5\,\%:
  \[
    \mathrm{CPI} = 1 + 0.02 \cdot 3 + 0.05 \cdot 5 = 1 + 0.06 + 0.25
    = 1.31 .
  \]
\end{example}

\begin{marginfigure}
  \resizebox{\linewidth}{!}{% A CPI stack (Lesson 8): the CPI of 1.31 of the five-stage example
% split into its base (steady-state) component of 1.00 and the
% contributions of branch mispredictions (0.06) and cache misses (0.25).
% Original.
% Usage: \resizebox{\linewidth}{!}{% A CPI stack (Lesson 8): the CPI of 1.31 of the five-stage example
% split into its base (steady-state) component of 1.00 and the
% contributions of branch mispredictions (0.06) and cache misses (0.25).
% Original.
% Usage: \resizebox{\linewidth}{!}{% A CPI stack (Lesson 8): the CPI of 1.31 of the five-stage example
% split into its base (steady-state) component of 1.00 and the
% contributions of branch mispredictions (0.06) and cache misses (0.25).
% Original.
% Usage: \resizebox{\linewidth}{!}{\input{figures/l08-cpi-stack}}
% Width: margin (50 mm); the drawing is about 45 mm wide.
\begin{tikzpicture}[x=1mm, y=24mm, font=\scriptsize\sffamily]
  \draw[fill=ln-accent!20] (0,0) rectangle (12,1.00)
    node[midway, align=center] {\iflangja{基本}{base}\\1.00};
  \draw[fill=ln-tip!30] (0,1.00) rectangle (12,1.06);
  \draw[fill=ln-caution!25] (0,1.06) rectangle (12,1.31);
  \draw[ln-muted] (12,1.03) -- (15,1.03)
    node[right, text=black] {\iflangja{分岐予測ミス 0.06}{branch mispred.\ 0.06}};
  \draw[ln-muted] (12,1.185) -- (15,1.185)
    node[right, text=black] {\iflangja{キャッシュミス 0.25}{cache misses 0.25}};
  \draw[->] (-2,0) -- (-2,1.45) node[above, inner sep=1pt] {CPI};
  \foreach \v in {0,0.5,1} {
    \draw (-2,\v) -- (-2.8,\v) node[left, inner sep=1pt] {\v};
  }
  \draw[dashed, ln-muted] (-2,1.31) -- (0,1.31);
  \node[left, inner sep=1pt] at (-2.8,1.31) {1.31};
\end{tikzpicture}
}
% Width: margin (50 mm); the drawing is about 45 mm wide.
\begin{tikzpicture}[x=1mm, y=24mm, font=\scriptsize\sffamily]
  \draw[fill=ln-accent!20] (0,0) rectangle (12,1.00)
    node[midway, align=center] {\iflangja{基本}{base}\\1.00};
  \draw[fill=ln-tip!30] (0,1.00) rectangle (12,1.06);
  \draw[fill=ln-caution!25] (0,1.06) rectangle (12,1.31);
  \draw[ln-muted] (12,1.03) -- (15,1.03)
    node[right, text=black] {\iflangja{分岐予測ミス 0.06}{branch mispred.\ 0.06}};
  \draw[ln-muted] (12,1.185) -- (15,1.185)
    node[right, text=black] {\iflangja{キャッシュミス 0.25}{cache misses 0.25}};
  \draw[->] (-2,0) -- (-2,1.45) node[above, inner sep=1pt] {CPI};
  \foreach \v in {0,0.5,1} {
    \draw (-2,\v) -- (-2.8,\v) node[left, inner sep=1pt] {\v};
  }
  \draw[dashed, ln-muted] (-2,1.31) -- (0,1.31);
  \node[left, inner sep=1pt] at (-2.8,1.31) {1.31};
\end{tikzpicture}
}
% Width: margin (50 mm); the drawing is about 45 mm wide.
\begin{tikzpicture}[x=1mm, y=24mm, font=\scriptsize\sffamily]
  \draw[fill=ln-accent!20] (0,0) rectangle (12,1.00)
    node[midway, align=center] {\iflangja{基本}{base}\\1.00};
  \draw[fill=ln-tip!30] (0,1.00) rectangle (12,1.06);
  \draw[fill=ln-caution!25] (0,1.06) rectangle (12,1.31);
  \draw[ln-muted] (12,1.03) -- (15,1.03)
    node[right, text=black] {\iflangja{分岐予測ミス 0.06}{branch mispred.\ 0.06}};
  \draw[ln-muted] (12,1.185) -- (15,1.185)
    node[right, text=black] {\iflangja{キャッシュミス 0.25}{cache misses 0.25}};
  \draw[->] (-2,0) -- (-2,1.45) node[above, inner sep=1pt] {CPI};
  \foreach \v in {0,0.5,1} {
    \draw (-2,\v) -- (-2.8,\v) node[left, inner sep=1pt] {\v};
  }
  \draw[dashed, ln-muted] (-2,1.31) -- (0,1.31);
  \node[left, inner sep=1pt] at (-2.8,1.31) {1.31};
\end{tikzpicture}
}
  \caption{CPI stack of \cref{ex:l08-cpi}.}
  \label{fig:l08-cpi-stack}
\end{marginfigure}

The terms of the sum are worth showing separately, because they say
where the time goes.

\begin{definition}[CPI stack]
  \label{def:l08-cpi-stack}
  A \term{CPI stack}[CPIスタック] is a stacked bar of the components of
  the CPI: the base (steady-state) component and one component per
  kind of event (\cref{fig:l08-cpi-stack}).
\end{definition}

In the example, cache misses account for $0.25/1.31 \approx 19\,\%$ of
the time and mispredictions for under 5\,\%: a better cache helps more
than a better predictor.\tip{IPC is the reciprocal of CPI: a CPI of
1.31 is an IPC of $1/1.31 \approx 0.76$.}

A multithreaded processor overlaps the stalls of one thread with the
work of others.  Ideally, with $W_{\text{depth}}$ threads, the stalls
are divided among them:
\begin{equation}
  \mathrm{CPI}_{\text{MT}} \approx
    \frac{\mathrm{CPI}_{\text{single}}}{W_{\text{depth}}}
    + \text{contention terms}.
  \label{eq:l08-cpi-mt}
\end{equation}
The contention terms account for shared resources that run out when
many threads are active: busy execution units, DRAM bandwidth, and the
registers that track outstanding misses.

\begin{definition}[MSHR]
  \label{def:l08-mshr}
  An \term{MSHR (miss status holding register)}[MSHR]\index{miss status holding register|see{MSHR (miss status holding register)}} is an entry
  that records one outstanding cache miss (the missing block and the
  requests waiting for it), so that the cache can keep serving other
  accesses while the miss is serviced.  The number of MSHRs bounds the
  number of misses in flight.
\end{definition}
\margin{The MSHRs therefore cap the memory-level parallelism of
Lesson~4: with, for example, 32 MSHRs (an illustrative number), a 33rd outstanding miss stalls until one
returns, however many warps are ready.}

\begin{note}
  \label{note:l08-cpi-mt}
  Read \eqref{eq:l08-cpi-mt} as cycles per instruction of the whole
  processor, counting the instructions of all threads.  The ideal term
  cannot fall below the issue limit (1 for a pipeline that issues one
  instruction per cycle): with the CPI of 1.31 of \cref{ex:l08-cpi},
  two threads already reach that limit.  A GPU is a multithreaded
  processor with many warps, so the CPI models apply to it, but the
  contention terms, especially MSHRs and DRAM bandwidth, dominate.
\end{note}

\subsection{Interval analysis}

The CPI model treats every penalty as a fixed number of cycles added to
an otherwise perfect pipeline.  Interval analysis looks at the shape of
the execution over time.

\begin{definition}[Interval analysis]
  \label{def:l08-interval}
  The method of \term{interval analysis}[インターバル解析] divides an
  execution into intervals.  Each interval consists of a steady state, in which
  instructions are dispatched at the \emph{effective dispatch rate},
  followed by one \term{miss event}[ミスイベント], a long-latency event
  (a branch misprediction, a cache or TLB miss) during whose penalty
  no instruction is dispatched (\cref{fig:l08-interval}).  The
  execution time is the sum of the steady-state times and the
  penalties.
\end{definition}

\begin{figure}[tb]
  \centering
  % Interval analysis (Lesson 8): dispatch rate over time.  Each interval
% is a steady state at the effective dispatch rate followed by one miss
% event.  Dashed: the simple model, in which dispatch stops abruptly for
% the penalty.  Solid: the first-order model, with the pipeline drain
% before and the refill after the penalty.
% Redraws the outline figure m8l3_interval_based (both panels merged).
% Usage: % Interval analysis (Lesson 8): dispatch rate over time.  Each interval
% is a steady state at the effective dispatch rate followed by one miss
% event.  Dashed: the simple model, in which dispatch stops abruptly for
% the penalty.  Solid: the first-order model, with the pipeline drain
% before and the refill after the penalty.
% Redraws the outline figure m8l3_interval_based (both panels merged).
% Usage: % Interval analysis (Lesson 8): dispatch rate over time.  Each interval
% is a steady state at the effective dispatch rate followed by one miss
% event.  Dashed: the simple model, in which dispatch stops abruptly for
% the penalty.  Solid: the first-order model, with the pipeline drain
% before and the refill after the penalty.
% Redraws the outline figure m8l3_interval_based (both panels merged).
% Usage: \input{figures/l08-interval}
% Width: body (120 mm); the drawing is about 112 mm wide.
\begin{tikzpicture}[x=1mm, y=1mm, font=\scriptsize\sffamily,
    >={Stealth[length=1.5mm]}]
  \draw[->] (0,0) -- (110,0) node[below left, inner sep=1pt]
    {\iflangja{時間}{time}};
  \draw[->] (0,0) -- (0,24);
  \node[rotate=90, align=center, inner sep=1pt] at (-7,8)
    {\iflangja{ディスパッチ率}{dispatch rate}};
  \node[left, inner sep=1pt] at (0,14) {$D$};
  % simple model (dashed)
  \draw[dashed, ln-muted] (2,14) -- (36,14) -- (36,0.3) -- (56,0.3)
    -- (56,14) -- (90,14) -- (90,0.3) -- (104,0.3);
  % first-order model (solid)
  \draw[very thick, ln-accent] (2,14) -- (32,14) -- (38,0.6) -- (50,0.6)
    -- (58,14) -- (86,14) -- (92,0.6) -- (104,0.6);
  % labels
  \draw[->, ln-caution] (32,21) node[above, text=ln-caution, inner sep=1pt]
    {\iflangja{ミスイベント}{miss event}} -- (32,15);
  \draw[->, ln-caution] (86,21) node[above, text=ln-caution, inner sep=1pt]
    {\iflangja{ミスイベント}{miss event}} -- (86,15);
  \node[above, text=ln-accent] at (17,14)
    {\iflangja{定常状態}{steady state}};
  \node[text=ln-muted] at (35,6) [left] {\iflangja{ドレイン}{drain}};
  \node[text=ln-muted] at (55,6) [right] {\iflangja{リフィル}{refill}};
  \node[above, inner sep=1pt] at (45,1) {\iflangja{ペナルティ}{penalty}};
  % intervals
  \draw[|-|] (2,-4) -- (56,-4) node[midway, below, inner sep=1.5pt]
    {\iflangja{インターバル 1}{interval 1}};
  \draw[|-|] (56,-4) -- (104,-4) node[midway, below, inner sep=1.5pt]
    {\iflangja{インターバル 2}{interval 2}};
\end{tikzpicture}

% Width: body (120 mm); the drawing is about 112 mm wide.
\begin{tikzpicture}[x=1mm, y=1mm, font=\scriptsize\sffamily,
    >={Stealth[length=1.5mm]}]
  \draw[->] (0,0) -- (110,0) node[below left, inner sep=1pt]
    {\iflangja{時間}{time}};
  \draw[->] (0,0) -- (0,24);
  \node[rotate=90, align=center, inner sep=1pt] at (-7,8)
    {\iflangja{ディスパッチ率}{dispatch rate}};
  \node[left, inner sep=1pt] at (0,14) {$D$};
  % simple model (dashed)
  \draw[dashed, ln-muted] (2,14) -- (36,14) -- (36,0.3) -- (56,0.3)
    -- (56,14) -- (90,14) -- (90,0.3) -- (104,0.3);
  % first-order model (solid)
  \draw[very thick, ln-accent] (2,14) -- (32,14) -- (38,0.6) -- (50,0.6)
    -- (58,14) -- (86,14) -- (92,0.6) -- (104,0.6);
  % labels
  \draw[->, ln-caution] (32,21) node[above, text=ln-caution, inner sep=1pt]
    {\iflangja{ミスイベント}{miss event}} -- (32,15);
  \draw[->, ln-caution] (86,21) node[above, text=ln-caution, inner sep=1pt]
    {\iflangja{ミスイベント}{miss event}} -- (86,15);
  \node[above, text=ln-accent] at (17,14)
    {\iflangja{定常状態}{steady state}};
  \node[text=ln-muted] at (35,6) [left] {\iflangja{ドレイン}{drain}};
  \node[text=ln-muted] at (55,6) [right] {\iflangja{リフィル}{refill}};
  \node[above, inner sep=1pt] at (45,1) {\iflangja{ペナルティ}{penalty}};
  % intervals
  \draw[|-|] (2,-4) -- (56,-4) node[midway, below, inner sep=1.5pt]
    {\iflangja{インターバル 1}{interval 1}};
  \draw[|-|] (56,-4) -- (104,-4) node[midway, below, inner sep=1.5pt]
    {\iflangja{インターバル 2}{interval 2}};
\end{tikzpicture}

% Width: body (120 mm); the drawing is about 112 mm wide.
\begin{tikzpicture}[x=1mm, y=1mm, font=\scriptsize\sffamily,
    >={Stealth[length=1.5mm]}]
  \draw[->] (0,0) -- (110,0) node[below left, inner sep=1pt]
    {\iflangja{時間}{time}};
  \draw[->] (0,0) -- (0,24);
  \node[rotate=90, align=center, inner sep=1pt] at (-7,8)
    {\iflangja{ディスパッチ率}{dispatch rate}};
  \node[left, inner sep=1pt] at (0,14) {$D$};
  % simple model (dashed)
  \draw[dashed, ln-muted] (2,14) -- (36,14) -- (36,0.3) -- (56,0.3)
    -- (56,14) -- (90,14) -- (90,0.3) -- (104,0.3);
  % first-order model (solid)
  \draw[very thick, ln-accent] (2,14) -- (32,14) -- (38,0.6) -- (50,0.6)
    -- (58,14) -- (86,14) -- (92,0.6) -- (104,0.6);
  % labels
  \draw[->, ln-caution] (32,21) node[above, text=ln-caution, inner sep=1pt]
    {\iflangja{ミスイベント}{miss event}} -- (32,15);
  \draw[->, ln-caution] (86,21) node[above, text=ln-caution, inner sep=1pt]
    {\iflangja{ミスイベント}{miss event}} -- (86,15);
  \node[above, text=ln-accent] at (17,14)
    {\iflangja{定常状態}{steady state}};
  \node[text=ln-muted] at (35,6) [left] {\iflangja{ドレイン}{drain}};
  \node[text=ln-muted] at (55,6) [right] {\iflangja{リフィル}{refill}};
  \node[above, inner sep=1pt] at (45,1) {\iflangja{ペナルティ}{penalty}};
  % intervals
  \draw[|-|] (2,-4) -- (56,-4) node[midway, below, inner sep=1.5pt]
    {\iflangja{インターバル 1}{interval 1}};
  \draw[|-|] (56,-4) -- (104,-4) node[midway, below, inner sep=1.5pt]
    {\iflangja{インターバル 2}{interval 2}};
\end{tikzpicture}

  \caption{Interval analysis.  Each interval is a steady state at the
    effective dispatch rate $D$ followed by one miss event.  Dashed: the
    abrupt stop of the simple model.  Solid: the first-order model with
    the pipeline drain before and the refill after the penalty.}
  \label{fig:l08-interval}
\end{figure}

A pipeline does not stop and restart abruptly.  When a miss event
occurs, instructions already in flight still drain out of the pipeline;
after it is resolved, the pipeline needs time to fill again before
dispatch reaches the steady rate.  Karkhanis and Smith's first-order
superscalar model adds these \emph{drain} and \emph{refill} phases,
whose shapes depend on the kind of event: a misprediction is detected
only when the branch executes, while a long cache miss blocks the
head of the reorder buffer and stalls dispatch once it fills
\cite{karkhanis2004firstorder}.\margin{Eeckhout \cite{eeckhout2010perfeval} presents interval analysis and the
first-order model in detail, including the penalties of each kind of
miss event.}

Applied to a GPU, the naive approach treats the GPU as one more
multithreaded processor.  That misses the two effects that most
distinguish GPUs:
\begin{description}
  \item[Branch divergence.]  Not every lane of a warp is active, and
    the paths of a divergent branch are serialised (Lesson~6), so a
    warp instruction does not represent a fixed amount of work.
  \item[Memory divergence.]  The latency of a warp's memory
    instruction depends on whether its accesses coalesce: an
    uncoalesced access generates many transactions and waits for the
    last of them (Lesson~4).
\end{description}
\margin{GPUMech applies interval analysis to GPUs, models the
round-robin and greedy-then-oldest schedulers of Lesson~7 with 13.2\,\%
and 14.0\,\% error at 97 times the speed of simulation, and produces
CPI stacks \cite{huang2014gpumech}.}%
More recent GPU interval models add sub-cores, sectored caches and the
contention for other shared resources.

\begin{note}[A GPU model of memory and computation parallelism]
  \label{note:l08-mwp-cwp}
  Hong and Kim's analytical model of a GPU, one of this unit's optional
  readings, rests on two quantities \cite{hong2009analytical}.
  \term{memory warp parallelism (MWP)}[メモリワープ並列度]\index{MWP|see{memory warp parallelism (MWP)}} is the
  maximum number of warps per SM that can access memory at the same
  time, during one warp's memory waiting period; it is limited by the
  latency, the departure delay between consecutive requests, the
  bandwidth and the number of active warps $N$:
  \[
    \mathrm{MWP} = \min(\mathrm{MWP}_{\text{without BW}},
    \mathrm{MWP}_{\text{peak BW}}, N),
  \]
  where
  $\mathrm{MWP}_{\text{without BW}} =
  \textit{Mem\_L}/\textit{Departure\_delay}$ and
  $\mathrm{MWP}_{\text{peak BW}}$ is the peak bandwidth divided by the
  bandwidth one warp needs times the number of active SMs.
  \term{computation warp parallelism (CWP)}[計算ワープ並列度]\index{CWP|see{computation warp parallelism (CWP)}} is the
  number of warps that the SM can execute during one memory waiting
  period, plus one:
  $\mathrm{CWP} = \min\bigl((\textit{Mem\_cycles} +
  \textit{Comp\_cycles})/\textit{Comp\_cycles},\; N\bigr)$.  If
  $\mathrm{CWP} \geq \mathrm{MWP}$ the computation of the other warps
  covers more than the memory system can overlap: the kernel is
  memory-bound and its time grows with $N/\mathrm{MWP}$ memory
  periods.  If $\mathrm{MWP} > \mathrm{CWP}$ the memory latency is
  hidden and the time is about one memory latency plus
  $N \times \textit{Comp\_cycles}$.
\end{note}
\margin{Hong and Kim report a geometric-mean error of 5.4\,\% on
micro-benchmarks and 13.3\,\% on applications
\cite{hong2009analytical}.}

\begin{keypoint}
  Analytical models trade accuracy for insight.  Roofline compares a
  kernel's arithmetic intensity with the machine's ridge point to decide
  whether it is compute-bound or memory-bound.  CPI models add
  frequency-weighted penalties to a steady-state CPI, and a CPI stack
  shows where the time goes; interval analysis refines this with drain
  and refill, and GPU versions must add branch and memory divergence.
\end{keypoint}

\section{Accelerating GPU simulation}
\label{sec:l08-accelerating}

Detailed simulation is slow.\source{M8, accelerating simulation;
Eeckhout \cite{eeckhout2010perfeval} (sampling); \cite{liu2023photon,sherwood2002simpoint}.}
A cycle-level simulator that models every unit in every cycle runs at
the order of a thousand simulated instructions per second (\num{1}
KIPS).  At that speed the \num{e9} instructions that a processor
executes in about a second take about 28~hours to simulate, and a
machine-learning workload that runs for hours on real hardware would
take on the order of a century.\caution{Lecture vs arithmetic:
$\num{e9}/\num{e3} = \num{e6}$~s, about 11.6~days, not 28~hours; the
28~hours correspond to 10~KIPS ($\num{e5}$~s).  At 1~KIPS one hour of
native execution takes about $3600 \times 11.6~\text{days} \approx
114$~years, which is the century.}  There are two ways out: make the
simulator faster, or give it less to simulate.

\subsection{Faster simulators}

\paragraph{Parallel simulation.}  A simulator can run on several host
threads, for example one thread per simulated core.\margin{Speeds of
CPU simulators: gem5, Flexus and MARSS reach about 200~KIPS with
detailed models; zsim simulates a detailed out-of-order core at
20~MIPS and a 1024-core chip at up to 300~MIPS on a 16-core host
\cite{sanchez2013zsim}.}  The cores are
nearly independent, but the memory system and the network on chip are
shared: the threads must communicate, or serialise their accesses to
the shared model, at every interaction.  Most workloads are
memory-bound, so these interactions are frequent and the memory model
becomes the bottleneck of the parallel simulator.

\paragraph{Event-driven simulation.}

\begin{definition}[Event-driven simulation]
  \label{def:l08-event-driven}
  In \term{event-driven simulation}[イベント駆動シミュレーション]
  the simulator keeps a queue of future events, each with a time stamp,
  ordered by time.  It repeatedly removes the earliest event, advances
  the simulated time directly to its time stamp, and processes it;
  processing may insert new events (\cref{fig:l08-event-queue}).
\end{definition}

\begin{figure}[tb]
  \centering
  % Event-driven simulation (Lesson 8).  Top: an event queue ordered by
% time stamp; processing the event at cycle 100 inserts a new event at
% cycle 250 in time order.  Bottom: simulated time jumps from one event
% to the next, skipping the cycles in between.
% Redraws the outline figure m8l4_esim with our own events and times.
% Usage: % Event-driven simulation (Lesson 8).  Top: an event queue ordered by
% time stamp; processing the event at cycle 100 inserts a new event at
% cycle 250 in time order.  Bottom: simulated time jumps from one event
% to the next, skipping the cycles in between.
% Redraws the outline figure m8l4_esim with our own events and times.
% Usage: % Event-driven simulation (Lesson 8).  Top: an event queue ordered by
% time stamp; processing the event at cycle 100 inserts a new event at
% cycle 250 in time order.  Bottom: simulated time jumps from one event
% to the next, skipping the cycles in between.
% Redraws the outline figure m8l4_esim with our own events and times.
% Usage: \input{figures/l08-event-queue}
% Width: body (120 mm); the drawing is about 112 mm wide.
\begin{tikzpicture}[x=1mm, y=1mm, font=\scriptsize\sffamily,
    >={Stealth[length=1.5mm]},
    ev/.style={draw, fill=ln-box, minimum width=24mm, minimum height=8mm,
               align=center, inner sep=1pt}]
  \node[anchor=west, font=\scriptsize\sffamily\bfseries] at (0,40)
    {\iflangja{イベントキュー（時刻順）}{event queue (ordered by time)}};
  \node[ev, fill=ln-accent!15] (e1) at (14,31)
    {$t=100$\\\iflangja{DRAM 応答}{DRAM reply}};
  \node[ev] (e2) at (42,31)
    {$t=104$\\\iflangja{ワープ3 起床}{warp 3 wakes}};
  \node[ev] (e3) at (98,31)
    {$t=400$\\\iflangja{I/O 完了}{I/O done}};
  \node[ev, fill=ln-tip!15] (e4) at (70,31)
    {$t=250$\\\iflangja{L2 書き戻し}{L2 write-back}};
  \draw[->] (e1) -- (e2);
  \draw[->] (e2) -- (e4);
  \draw[->] (e4) -- (e3);
  \draw[->, ln-tip] (e1.south) to[bend right=18]
    node[pos=0.55, below, text=ln-tip, inner sep=1.5pt]
    {\iflangja{処理中に新しいイベントを挿入}{processing it inserts a new event}}
    (e4.south);
  % timeline
  \draw[->] (0,6) -- (110,6);
  \node[anchor=north east, inner sep=1pt] at (110,0.5)
    {\iflangja{シミュレーション時刻（サイクル）}{simulated time (cycles)}};
  \foreach \t/\x in {100/14,104/22,250/62,400/98} {
    \fill (\x,6) circle (0.8);
    \node[below, inner sep=2pt] at (\x,5) {\t};
  }
  \draw[->, ln-caution] (14,7.5) to[bend left=50] (22,7.5);
  \draw[->, ln-caution] (22,7.5) to[bend left=30]
    node[midway, above, text=ln-caution, inner sep=1pt]
    {\iflangja{145サイクルをスキップ}{skip 145 cycles}} (62,7.5);
  \draw[->, ln-caution] (62,7.5) to[bend left=30]
    node[midway, above, text=ln-caution, inner sep=1pt]
    {\iflangja{149サイクルをスキップ}{skip 149 cycles}} (98,7.5);
\end{tikzpicture}

% Width: body (120 mm); the drawing is about 112 mm wide.
\begin{tikzpicture}[x=1mm, y=1mm, font=\scriptsize\sffamily,
    >={Stealth[length=1.5mm]},
    ev/.style={draw, fill=ln-box, minimum width=24mm, minimum height=8mm,
               align=center, inner sep=1pt}]
  \node[anchor=west, font=\scriptsize\sffamily\bfseries] at (0,40)
    {\iflangja{イベントキュー（時刻順）}{event queue (ordered by time)}};
  \node[ev, fill=ln-accent!15] (e1) at (14,31)
    {$t=100$\\\iflangja{DRAM 応答}{DRAM reply}};
  \node[ev] (e2) at (42,31)
    {$t=104$\\\iflangja{ワープ3 起床}{warp 3 wakes}};
  \node[ev] (e3) at (98,31)
    {$t=400$\\\iflangja{I/O 完了}{I/O done}};
  \node[ev, fill=ln-tip!15] (e4) at (70,31)
    {$t=250$\\\iflangja{L2 書き戻し}{L2 write-back}};
  \draw[->] (e1) -- (e2);
  \draw[->] (e2) -- (e4);
  \draw[->] (e4) -- (e3);
  \draw[->, ln-tip] (e1.south) to[bend right=18]
    node[pos=0.55, below, text=ln-tip, inner sep=1.5pt]
    {\iflangja{処理中に新しいイベントを挿入}{processing it inserts a new event}}
    (e4.south);
  % timeline
  \draw[->] (0,6) -- (110,6);
  \node[anchor=north east, inner sep=1pt] at (110,0.5)
    {\iflangja{シミュレーション時刻（サイクル）}{simulated time (cycles)}};
  \foreach \t/\x in {100/14,104/22,250/62,400/98} {
    \fill (\x,6) circle (0.8);
    \node[below, inner sep=2pt] at (\x,5) {\t};
  }
  \draw[->, ln-caution] (14,7.5) to[bend left=50] (22,7.5);
  \draw[->, ln-caution] (22,7.5) to[bend left=30]
    node[midway, above, text=ln-caution, inner sep=1pt]
    {\iflangja{145サイクルをスキップ}{skip 145 cycles}} (62,7.5);
  \draw[->, ln-caution] (62,7.5) to[bend left=30]
    node[midway, above, text=ln-caution, inner sep=1pt]
    {\iflangja{149サイクルをスキップ}{skip 149 cycles}} (98,7.5);
\end{tikzpicture}

% Width: body (120 mm); the drawing is about 112 mm wide.
\begin{tikzpicture}[x=1mm, y=1mm, font=\scriptsize\sffamily,
    >={Stealth[length=1.5mm]},
    ev/.style={draw, fill=ln-box, minimum width=24mm, minimum height=8mm,
               align=center, inner sep=1pt}]
  \node[anchor=west, font=\scriptsize\sffamily\bfseries] at (0,40)
    {\iflangja{イベントキュー（時刻順）}{event queue (ordered by time)}};
  \node[ev, fill=ln-accent!15] (e1) at (14,31)
    {$t=100$\\\iflangja{DRAM 応答}{DRAM reply}};
  \node[ev] (e2) at (42,31)
    {$t=104$\\\iflangja{ワープ3 起床}{warp 3 wakes}};
  \node[ev] (e3) at (98,31)
    {$t=400$\\\iflangja{I/O 完了}{I/O done}};
  \node[ev, fill=ln-tip!15] (e4) at (70,31)
    {$t=250$\\\iflangja{L2 書き戻し}{L2 write-back}};
  \draw[->] (e1) -- (e2);
  \draw[->] (e2) -- (e4);
  \draw[->] (e4) -- (e3);
  \draw[->, ln-tip] (e1.south) to[bend right=18]
    node[pos=0.55, below, text=ln-tip, inner sep=1.5pt]
    {\iflangja{処理中に新しいイベントを挿入}{processing it inserts a new event}}
    (e4.south);
  % timeline
  \draw[->] (0,6) -- (110,6);
  \node[anchor=north east, inner sep=1pt] at (110,0.5)
    {\iflangja{シミュレーション時刻（サイクル）}{simulated time (cycles)}};
  \foreach \t/\x in {100/14,104/22,250/62,400/98} {
    \fill (\x,6) circle (0.8);
    \node[below, inner sep=2pt] at (\x,5) {\t};
  }
  \draw[->, ln-caution] (14,7.5) to[bend left=50] (22,7.5);
  \draw[->, ln-caution] (22,7.5) to[bend left=30]
    node[midway, above, text=ln-caution, inner sep=1pt]
    {\iflangja{145サイクルをスキップ}{skip 145 cycles}} (62,7.5);
  \draw[->, ln-caution] (62,7.5) to[bend left=30]
    node[midway, above, text=ln-caution, inner sep=1pt]
    {\iflangja{149サイクルをスキップ}{skip 149 cycles}} (98,7.5);
\end{tikzpicture}

  \caption{Event-driven simulation.  The event queue is ordered by time
    stamp; processing an event can insert a new one.  Simulated time
    jumps from event to event and skips the cycles in between.}
  \label{fig:l08-event-queue}
\end{figure}

The cycles between events cost nothing, so event-driven simulation
pays off when events are sparse, as with long latencies: an I/O
operation or a write to non-volatile memory that takes thousands of
cycles is one event instead of thousands of ticks.\tip{Cycle-level
simulation costs time per cycle, event-driven simulation per
event.  A busy pipeline has events in nearly every cycle, so its core
gains little; long-latency devices gain most.}

\paragraph{Simpler models.}  Not every event matters equally.  A
simulator can use average latencies or throughputs for non-critical
events and keep detail only where the question requires it.  To study
the memory hierarchy, the core can be reduced to issuing instructions
at an IPC equal to its issue width, with only the caches and memory
modeled; zsim, a fast simulator of many-core CPUs, offers such a
simple core (one instruction per cycle) next to its detailed
out-of-order model \cite{sanchez2013zsim}.  To study the core,
the memory can be given a fixed latency.\caution{Every simplification
removes an interaction.  A fixed memory latency, for example, hides the
queueing delays that a bandwidth-hungry design creates, exactly the
contention terms of \eqref{eq:l08-cpi-mt}.}

The remaining two accelerations, sampling and statistical models, are
the subject of the rest of this section.

\subsection{Sampling}

\begin{definition}[Sampled simulation]
  \label{def:l08-sampled}
  In \term{sampled simulation}[サンプリングシミュレーション] only
  selected parts of a program's execution, the samples, are simulated
  in detail, and the performance of the whole is estimated from them.
  The samples may be chosen at random or by analysing the program.
\end{definition}

The parts between samples still have to be skipped somehow.  An
execution-driven simulator fast-forwards through them by executing
functionally, without timing, which is much faster, or starts each
sample from a \term{checkpoint}[チェックポイント], a saved
architectural state (registers and memory) from which the simulation
can resume without executing what came before.  A trace-driven
simulator simply collects traces of the samples only.

A sample does not start from a realistic microarchitectural state:
its caches and branch predictors are empty or hold the contents left
by the previous sample.  The \term{warm-up}[ウォームアップ] of a sample
is a stretch of execution simulated before it, whose timing is
discarded, during which these structures are filled.  Alternatively
the samples are made long enough for the cold start to be negligible.

Random samples need many draws for an accurate estimate.  Programs,
however, are not random: they repeat the same behaviour.

\begin{definition}[Program phase and SimPoint]
  \label{def:l08-phase}
  A \term{program phase}[プログラムフェーズ] is a stretch of execution
  with homogeneous behaviour (IPC, miss rates), which may recur many
  times.  \term{SimPoint}[SimPoint] identifies the phases of a
  program by clustering intervals of its execution by the code they
  execute, and selects one representative interval per phase, weighted
  by the fraction of execution the phase accounts for
  (\cref{fig:l08-phases}) \cite{sherwood2002simpoint}.  Each interval
  is described by a \emph{basic block vector}, how often each basic
  block executed in it, and the vectors are grouped by $k$-means
  clustering.
\end{definition}

\begin{marginfigure}
  \resizebox{\linewidth}{!}{% Program phases and phase-based sampling (Lesson 8): IPC over time for
% a program that alternates between a high-IPC phase A and a low-IPC
% phase B.  Phase-based sampling (SimPoint) simulates one representative
% interval per phase (outlined) and weights it by the phase's share.
% Redraws the outline figure m8l4_sample with our own phase lengths.
% Usage: \resizebox{\linewidth}{!}{% Program phases and phase-based sampling (Lesson 8): IPC over time for
% a program that alternates between a high-IPC phase A and a low-IPC
% phase B.  Phase-based sampling (SimPoint) simulates one representative
% interval per phase (outlined) and weights it by the phase's share.
% Redraws the outline figure m8l4_sample with our own phase lengths.
% Usage: \resizebox{\linewidth}{!}{% Program phases and phase-based sampling (Lesson 8): IPC over time for
% a program that alternates between a high-IPC phase A and a low-IPC
% phase B.  Phase-based sampling (SimPoint) simulates one representative
% interval per phase (outlined) and weights it by the phase's share.
% Redraws the outline figure m8l4_sample with our own phase lengths.
% Usage: \resizebox{\linewidth}{!}{\input{figures/l08-phases}}
% Width: margin (50 mm); the drawing is about 50 mm wide.
\begin{tikzpicture}[x=1.15mm, y=1mm, font=\scriptsize\sffamily]
  \foreach \a/\b/\p/\h/\c in {0/8/A/16/ln-tip!30,8/14/B/6/ln-source!25,
      14/22/A/16/ln-tip!30,22/28/B/6/ln-source!25,28/36/A/16/ln-tip!30,
      36/40/B/6/ln-source!25} {
    \fill[\c] (\a,0) rectangle (\b,\h);
    \node at ({(\a+\b)/2},-2.2) {\p};
  }
  \draw[very thick, ln-caution] (2,0) rectangle (5,16);
  \draw[very thick, ln-caution] (9.5,0) rectangle (12.5,6);
  \draw[->] (0,0) -- (41.5,0);
  \draw[->] (0,0) -- (0,19) node[above right, inner sep=1pt] {IPC};
  \node[anchor=north east] at (41.5,-3.8) {\iflangja{時間}{time}};
\end{tikzpicture}
}
% Width: margin (50 mm); the drawing is about 50 mm wide.
\begin{tikzpicture}[x=1.15mm, y=1mm, font=\scriptsize\sffamily]
  \foreach \a/\b/\p/\h/\c in {0/8/A/16/ln-tip!30,8/14/B/6/ln-source!25,
      14/22/A/16/ln-tip!30,22/28/B/6/ln-source!25,28/36/A/16/ln-tip!30,
      36/40/B/6/ln-source!25} {
    \fill[\c] (\a,0) rectangle (\b,\h);
    \node at ({(\a+\b)/2},-2.2) {\p};
  }
  \draw[very thick, ln-caution] (2,0) rectangle (5,16);
  \draw[very thick, ln-caution] (9.5,0) rectangle (12.5,6);
  \draw[->] (0,0) -- (41.5,0);
  \draw[->] (0,0) -- (0,19) node[above right, inner sep=1pt] {IPC};
  \node[anchor=north east] at (41.5,-3.8) {\iflangja{時間}{time}};
\end{tikzpicture}
}
% Width: margin (50 mm); the drawing is about 50 mm wide.
\begin{tikzpicture}[x=1.15mm, y=1mm, font=\scriptsize\sffamily]
  \foreach \a/\b/\p/\h/\c in {0/8/A/16/ln-tip!30,8/14/B/6/ln-source!25,
      14/22/A/16/ln-tip!30,22/28/B/6/ln-source!25,28/36/A/16/ln-tip!30,
      36/40/B/6/ln-source!25} {
    \fill[\c] (\a,0) rectangle (\b,\h);
    \node at ({(\a+\b)/2},-2.2) {\p};
  }
  \draw[very thick, ln-caution] (2,0) rectangle (5,16);
  \draw[very thick, ln-caution] (9.5,0) rectangle (12.5,6);
  \draw[->] (0,0) -- (41.5,0);
  \draw[->] (0,0) -- (0,19) node[above right, inner sep=1pt] {IPC};
  \node[anchor=north east] at (41.5,-3.8) {\iflangja{時間}{time}};
\end{tikzpicture}
}
  \caption{IPC over time of a program alternating a high-IPC phase A
    and a low-IPC phase B.  Phase-based sampling simulates one
    representative of each (outlined).}
  \label{fig:l08-phases}
\end{marginfigure}

With a few phases recurring hundreds of times, simulating one interval
of each phase and weighting it gives most of the accuracy of a full
simulation at a small fraction of its cost.

\subsection{Sampling on GPUs}

A single GPU thread is too short to sample within, but a GPU program
consists of very many threads executing the same code (SPMD, Lesson~1),
which offers other granularities:
\begin{description}
  \item[Block level.]  Simulate about 100 thread blocks instead of the
    thousands of the full grid.
  \item[Kernel level.]  Simulate one or two kernels instead of tens.
  \item[Warp level.]  Simulate fewer warps of a kernel in detail.
\end{description}

\begin{definition}[Kernel-level sampling]
  \label{def:l08-kernel-sampling}
  \term{kernel-level sampling}[カーネル単位サンプリング] simulates in
  detail one launch of each distinct kernel, or of each group of
  similar launches, and uses its time for the other launches.
\end{definition}

Kernel-level sampling is especially effective for machine-learning
workloads, which launch the same kernels hundreds of times, once per
layer and iteration.  Photon, one of this unit's assigned readings,
combines sampling at three levels, kernels, warps and basic blocks, and
finds the repetitive behaviour by a lightweight online analysis instead
of profiling the workload beforehand \cite{liu2023photon}.\margin{Photon
reduces the simulation of one ResNet-152 inference (batch size~1) from
7.05~days to 1.7~hours with a sampling error of 10.7\,\%; it is built
on MGPUSim \cite{liu2023photon}.}

\subsection{Reducing the workload}

The other way out is to simulate less work.
\begin{itemize}
  \item \emph{Fewer iterations}: one training iteration instead of a
    thousand, since the iterations of a machine-learning training run
    behave alike.
  \item \emph{Smaller inputs}: a graph of \num{e6} nodes instead of
    \num{e9}.
  \item \emph{Only the hot code}: when one or two of hundreds of
    functions take 90\,\% of the time, simulate only those.
\end{itemize}
\caution{A smaller input can change the behaviour, not only shorten
it: a working set that fits in the L2 cache at reduced size may not fit
at full size, so the reduced run can hide the memory-bound behaviour
of the real one.}

\begin{definition}[Microbenchmark and representative workload]
  \label{def:l08-microbenchmark}
  A \term{microbenchmark}[マイクロベンチマーク] is a small program
  written to exercise one feature of a machine (a cache level, a
  memory access pattern, an instruction's latency) in isolation.  A
  \term{representative workload}[代表ワークロード] is a reduced program
  or input built so that its behaviour matches that of the full
  workload on the features under study.
\end{definition}

\subsection{Data-driven models}

The last approach does not model the machine at all.

\begin{definition}[Data-driven model]
  \label{def:l08-data-driven}
  A \term{data-driven model}[データ駆動モデル] predicts performance
  from features of a program (instruction count, number of loads,
  number of divergent branches, \dots) by a statistical or
  machine-learning model fitted to data, without encoding how the
  hardware works.
\end{definition}

\Needspace{20\baselineskip}
The procedure has three steps: run a set of programs on the hardware
(or a detailed simulator) and record their performance, for example
their IPC; extract features from hardware performance counters or from
the simulator; and fit a model.\margin{The hardware counters are the
ones that Nsight Compute (Lesson~5) reads: instructions executed per
type, memory transactions, branch divergence.}  A simple form of such a model is
\begin{equation}
  \textit{cycles} = C_1 N_{\text{inst}}^{e_1}
    + C_2 N_{\text{FP}}^{e_2} + C_3 N_{\text{mem}}^{e_3},
  \label{eq:l08-data-driven}
\end{equation}
where $N_{\text{inst}}$, $N_{\text{FP}}$ and $N_{\text{mem}}$ are the
counts of all, floating-point and memory instructions and the
coefficients $C_k$ and exponents $e_k$ are fitted to the
data.\caution{An analytical model is built from an understanding of
the machine; a data-driven model only from data.  It is accurate for
programs like those it was fitted on and gives no guarantee, and
little insight, outside them.}

\begin{keypoint}[Lesson summary]
  Cycle-level simulators advance a clock and model every unit every
  cycle, driven by execution or by traces; the pipeline is a set of
  FIFO queues with depths and widths plus a scheduler, and a GPU
  simulator models warps, their masks or SIMT stack, coalescing and
  sectored caches.  Analytical models give first-order answers: the FMA
  units, bandwidth and $W_{\text{depth}} \approx L/n$ of a design, the
  Roofline bound $\min(P_{\text{peak}}, B_{\text{peak}} \cdot I)$ with
  its ridge point, CPI stacks and interval analysis with branch and
  memory divergence.  Simulation at about 1~KIPS is too slow for real
  workloads; parallel and event-driven simulation, simplified models,
  sampling (phases, SimPoint, block-, kernel- and warp-level), reduced
  workloads and data-driven models make it tractable.
\end{keypoint}

\end{document}
